% latex-report skill: shared preamble
% Usage: % latex-report skill: shared preamble
% Usage: % latex-report skill: shared preamble
% Usage: \input{preamble} immediately after \documentclass is NOT needed;
% this file declares the class itself. Start your document with:
%   \input{preamble}
%   \begin{document}

\documentclass[11pt,a4paper]{article}

\usepackage[T1]{fontenc}
\usepackage[utf8]{inputenc}
\usepackage{lmodern}
\usepackage[margin=2.4cm,top=2.2cm,bottom=2.4cm]{geometry}
\usepackage{booktabs}
\usepackage{array}
\usepackage{longtable}
\usepackage{xcolor}
\usepackage{tcolorbox}
\usepackage{pgfplots}
\usepackage{caption}
\usepackage{titlesec}
\usepackage{enumitem}
\usepackage{fancyhdr}
\usepackage{microtype}
\usepackage[hidelinks]{hyperref}

\pgfplotsset{compat=1.18}
\usetikzlibrary{patterns}
\tcbuselibrary{skins}

% ---------------------------------------------------------------- palette
% Chosen to stay separable in greyscale and under deuteranopia.
\definecolor{palA}{HTML}{1F4E79}   % deep blue
\definecolor{palB}{HTML}{C55A11}   % burnt orange
\definecolor{palC}{HTML}{548235}   % green
\definecolor{palD}{HTML}{7030A0}   % purple
\definecolor{palE}{HTML}{808080}   % grey
\definecolor{rule}{HTML}{B0B0B0}
\definecolor{panelbg}{HTML}{F2F6FA}
\definecolor{panelln}{HTML}{1F4E79}

% ---------------------------------------------------------------- headings
\titleformat{\section}{\Large\bfseries\color{palA}}{\thesection}{0.6em}{}
\titleformat{\subsection}{\large\bfseries}{\thesubsection}{0.6em}{}
\titlespacing*{\section}{0pt}{1.4em}{0.6em}
\titlespacing*{\subsection}{0pt}{1.1em}{0.4em}

\captionsetup{font=small,labelfont={bf,color=palA},skip=6pt}
\setlist{itemsep=2pt,topsep=4pt,parsep=0pt}

\pagestyle{fancy}
\fancyhf{}
\renewcommand{\headrulewidth}{0.4pt}
\fancyfoot[C]{\small\thepage}

% ---------------------------------------------------------------- macros
% Boxed lead panel for the TL;DR.
\newcommand{\tldr}[1]{%
  \begin{tcolorbox}[colback=panelbg,colframe=panelln,boxrule=1pt,
    left=10pt,right=10pt,top=8pt,bottom=8pt,arc=2pt]
  #1
  \end{tcolorbox}\vspace{4pt}}

% Callout for a caveat or correction.
\newcommand{\caveat}[1]{%
  \begin{tcolorbox}[colback=white,colframe=palB,boxrule=0.8pt,
    left=9pt,right=9pt,top=6pt,bottom=6pt,arc=1pt]
  \small #1
  \end{tcolorbox}\vspace{3pt}}

\newcommand{\code}[1]{\texttt{\small #1}}
\newcommand{\num}[1]{\textbf{#1}}

% Numeric column that aligns on the decimal point.
\newcolumntype{R}{>{\raggedleft\arraybackslash}p{1.6cm}}

% Shared bar-chart styling. Use inside tikzpicture.
\pgfplotsset{
  barbase/.style={
    ybar,
    width=\linewidth,
    height=5.6cm,
    ymin=0,
    bar width=13pt,
    ymajorgrids=true,
    grid style={rule,dashed,very thin},
    axis line style={rule},
    tick align=outside,
    tick style={rule},
    legend style={draw=rule,font=\small,at={(0.5,1.04)},anchor=south,legend columns=-1},
    label style={font=\small},
    tick label style={font=\small},
    nodes near coords,
    nodes near coords style={font=\scriptsize,color=black!70},
  },
  stackbase/.style={
    ybar stacked,
    width=\linewidth,
    height=5.6cm,
    ymin=0,
    bar width=22pt,
    ymajorgrids=true,
    grid style={rule,dashed,very thin},
    axis line style={rule},
    legend style={draw=rule,font=\small,at={(1.02,1)},anchor=north west},
    label style={font=\small},
    tick label style={font=\small},
  },
}
 immediately after \documentclass is NOT needed;
% this file declares the class itself. Start your document with:
%   % latex-report skill: shared preamble
% Usage: \input{preamble} immediately after \documentclass is NOT needed;
% this file declares the class itself. Start your document with:
%   \input{preamble}
%   \begin{document}

\documentclass[11pt,a4paper]{article}

\usepackage[T1]{fontenc}
\usepackage[utf8]{inputenc}
\usepackage{lmodern}
\usepackage[margin=2.4cm,top=2.2cm,bottom=2.4cm]{geometry}
\usepackage{booktabs}
\usepackage{array}
\usepackage{longtable}
\usepackage{xcolor}
\usepackage{tcolorbox}
\usepackage{pgfplots}
\usepackage{caption}
\usepackage{titlesec}
\usepackage{enumitem}
\usepackage{fancyhdr}
\usepackage{microtype}
\usepackage[hidelinks]{hyperref}

\pgfplotsset{compat=1.18}
\usetikzlibrary{patterns}
\tcbuselibrary{skins}

% ---------------------------------------------------------------- palette
% Chosen to stay separable in greyscale and under deuteranopia.
\definecolor{palA}{HTML}{1F4E79}   % deep blue
\definecolor{palB}{HTML}{C55A11}   % burnt orange
\definecolor{palC}{HTML}{548235}   % green
\definecolor{palD}{HTML}{7030A0}   % purple
\definecolor{palE}{HTML}{808080}   % grey
\definecolor{rule}{HTML}{B0B0B0}
\definecolor{panelbg}{HTML}{F2F6FA}
\definecolor{panelln}{HTML}{1F4E79}

% ---------------------------------------------------------------- headings
\titleformat{\section}{\Large\bfseries\color{palA}}{\thesection}{0.6em}{}
\titleformat{\subsection}{\large\bfseries}{\thesubsection}{0.6em}{}
\titlespacing*{\section}{0pt}{1.4em}{0.6em}
\titlespacing*{\subsection}{0pt}{1.1em}{0.4em}

\captionsetup{font=small,labelfont={bf,color=palA},skip=6pt}
\setlist{itemsep=2pt,topsep=4pt,parsep=0pt}

\pagestyle{fancy}
\fancyhf{}
\renewcommand{\headrulewidth}{0.4pt}
\fancyfoot[C]{\small\thepage}

% ---------------------------------------------------------------- macros
% Boxed lead panel for the TL;DR.
\newcommand{\tldr}[1]{%
  \begin{tcolorbox}[colback=panelbg,colframe=panelln,boxrule=1pt,
    left=10pt,right=10pt,top=8pt,bottom=8pt,arc=2pt]
  #1
  \end{tcolorbox}\vspace{4pt}}

% Callout for a caveat or correction.
\newcommand{\caveat}[1]{%
  \begin{tcolorbox}[colback=white,colframe=palB,boxrule=0.8pt,
    left=9pt,right=9pt,top=6pt,bottom=6pt,arc=1pt]
  \small #1
  \end{tcolorbox}\vspace{3pt}}

\newcommand{\code}[1]{\texttt{\small #1}}
\newcommand{\num}[1]{\textbf{#1}}

% Numeric column that aligns on the decimal point.
\newcolumntype{R}{>{\raggedleft\arraybackslash}p{1.6cm}}

% Shared bar-chart styling. Use inside tikzpicture.
\pgfplotsset{
  barbase/.style={
    ybar,
    width=\linewidth,
    height=5.6cm,
    ymin=0,
    bar width=13pt,
    ymajorgrids=true,
    grid style={rule,dashed,very thin},
    axis line style={rule},
    tick align=outside,
    tick style={rule},
    legend style={draw=rule,font=\small,at={(0.5,1.04)},anchor=south,legend columns=-1},
    label style={font=\small},
    tick label style={font=\small},
    nodes near coords,
    nodes near coords style={font=\scriptsize,color=black!70},
  },
  stackbase/.style={
    ybar stacked,
    width=\linewidth,
    height=5.6cm,
    ymin=0,
    bar width=22pt,
    ymajorgrids=true,
    grid style={rule,dashed,very thin},
    axis line style={rule},
    legend style={draw=rule,font=\small,at={(1.02,1)},anchor=north west},
    label style={font=\small},
    tick label style={font=\small},
  },
}

%   \begin{document}

\documentclass[11pt,a4paper]{article}

\usepackage[T1]{fontenc}
\usepackage[utf8]{inputenc}
\usepackage{lmodern}
\usepackage[margin=2.4cm,top=2.2cm,bottom=2.4cm]{geometry}
\usepackage{booktabs}
\usepackage{array}
\usepackage{longtable}
\usepackage{xcolor}
\usepackage{tcolorbox}
\usepackage{pgfplots}
\usepackage{caption}
\usepackage{titlesec}
\usepackage{enumitem}
\usepackage{fancyhdr}
\usepackage{microtype}
\usepackage[hidelinks]{hyperref}

\pgfplotsset{compat=1.18}
\usetikzlibrary{patterns}
\tcbuselibrary{skins}

% ---------------------------------------------------------------- palette
% Chosen to stay separable in greyscale and under deuteranopia.
\definecolor{palA}{HTML}{1F4E79}   % deep blue
\definecolor{palB}{HTML}{C55A11}   % burnt orange
\definecolor{palC}{HTML}{548235}   % green
\definecolor{palD}{HTML}{7030A0}   % purple
\definecolor{palE}{HTML}{808080}   % grey
\definecolor{rule}{HTML}{B0B0B0}
\definecolor{panelbg}{HTML}{F2F6FA}
\definecolor{panelln}{HTML}{1F4E79}

% ---------------------------------------------------------------- headings
\titleformat{\section}{\Large\bfseries\color{palA}}{\thesection}{0.6em}{}
\titleformat{\subsection}{\large\bfseries}{\thesubsection}{0.6em}{}
\titlespacing*{\section}{0pt}{1.4em}{0.6em}
\titlespacing*{\subsection}{0pt}{1.1em}{0.4em}

\captionsetup{font=small,labelfont={bf,color=palA},skip=6pt}
\setlist{itemsep=2pt,topsep=4pt,parsep=0pt}

\pagestyle{fancy}
\fancyhf{}
\renewcommand{\headrulewidth}{0.4pt}
\fancyfoot[C]{\small\thepage}

% ---------------------------------------------------------------- macros
% Boxed lead panel for the TL;DR.
\newcommand{\tldr}[1]{%
  \begin{tcolorbox}[colback=panelbg,colframe=panelln,boxrule=1pt,
    left=10pt,right=10pt,top=8pt,bottom=8pt,arc=2pt]
  #1
  \end{tcolorbox}\vspace{4pt}}

% Callout for a caveat or correction.
\newcommand{\caveat}[1]{%
  \begin{tcolorbox}[colback=white,colframe=palB,boxrule=0.8pt,
    left=9pt,right=9pt,top=6pt,bottom=6pt,arc=1pt]
  \small #1
  \end{tcolorbox}\vspace{3pt}}

\newcommand{\code}[1]{\texttt{\small #1}}
\newcommand{\num}[1]{\textbf{#1}}

% Numeric column that aligns on the decimal point.
\newcolumntype{R}{>{\raggedleft\arraybackslash}p{1.6cm}}

% Shared bar-chart styling. Use inside tikzpicture.
\pgfplotsset{
  barbase/.style={
    ybar,
    width=\linewidth,
    height=5.6cm,
    ymin=0,
    bar width=13pt,
    ymajorgrids=true,
    grid style={rule,dashed,very thin},
    axis line style={rule},
    tick align=outside,
    tick style={rule},
    legend style={draw=rule,font=\small,at={(0.5,1.04)},anchor=south,legend columns=-1},
    label style={font=\small},
    tick label style={font=\small},
    nodes near coords,
    nodes near coords style={font=\scriptsize,color=black!70},
  },
  stackbase/.style={
    ybar stacked,
    width=\linewidth,
    height=5.6cm,
    ymin=0,
    bar width=22pt,
    ymajorgrids=true,
    grid style={rule,dashed,very thin},
    axis line style={rule},
    legend style={draw=rule,font=\small,at={(1.02,1)},anchor=north west},
    label style={font=\small},
    tick label style={font=\small},
  },
}
 immediately after \documentclass is NOT needed;
% this file declares the class itself. Start your document with:
%   % latex-report skill: shared preamble
% Usage: % latex-report skill: shared preamble
% Usage: \input{preamble} immediately after \documentclass is NOT needed;
% this file declares the class itself. Start your document with:
%   \input{preamble}
%   \begin{document}

\documentclass[11pt,a4paper]{article}

\usepackage[T1]{fontenc}
\usepackage[utf8]{inputenc}
\usepackage{lmodern}
\usepackage[margin=2.4cm,top=2.2cm,bottom=2.4cm]{geometry}
\usepackage{booktabs}
\usepackage{array}
\usepackage{longtable}
\usepackage{xcolor}
\usepackage{tcolorbox}
\usepackage{pgfplots}
\usepackage{caption}
\usepackage{titlesec}
\usepackage{enumitem}
\usepackage{fancyhdr}
\usepackage{microtype}
\usepackage[hidelinks]{hyperref}

\pgfplotsset{compat=1.18}
\usetikzlibrary{patterns}
\tcbuselibrary{skins}

% ---------------------------------------------------------------- palette
% Chosen to stay separable in greyscale and under deuteranopia.
\definecolor{palA}{HTML}{1F4E79}   % deep blue
\definecolor{palB}{HTML}{C55A11}   % burnt orange
\definecolor{palC}{HTML}{548235}   % green
\definecolor{palD}{HTML}{7030A0}   % purple
\definecolor{palE}{HTML}{808080}   % grey
\definecolor{rule}{HTML}{B0B0B0}
\definecolor{panelbg}{HTML}{F2F6FA}
\definecolor{panelln}{HTML}{1F4E79}

% ---------------------------------------------------------------- headings
\titleformat{\section}{\Large\bfseries\color{palA}}{\thesection}{0.6em}{}
\titleformat{\subsection}{\large\bfseries}{\thesubsection}{0.6em}{}
\titlespacing*{\section}{0pt}{1.4em}{0.6em}
\titlespacing*{\subsection}{0pt}{1.1em}{0.4em}

\captionsetup{font=small,labelfont={bf,color=palA},skip=6pt}
\setlist{itemsep=2pt,topsep=4pt,parsep=0pt}

\pagestyle{fancy}
\fancyhf{}
\renewcommand{\headrulewidth}{0.4pt}
\fancyfoot[C]{\small\thepage}

% ---------------------------------------------------------------- macros
% Boxed lead panel for the TL;DR.
\newcommand{\tldr}[1]{%
  \begin{tcolorbox}[colback=panelbg,colframe=panelln,boxrule=1pt,
    left=10pt,right=10pt,top=8pt,bottom=8pt,arc=2pt]
  #1
  \end{tcolorbox}\vspace{4pt}}

% Callout for a caveat or correction.
\newcommand{\caveat}[1]{%
  \begin{tcolorbox}[colback=white,colframe=palB,boxrule=0.8pt,
    left=9pt,right=9pt,top=6pt,bottom=6pt,arc=1pt]
  \small #1
  \end{tcolorbox}\vspace{3pt}}

\newcommand{\code}[1]{\texttt{\small #1}}
\newcommand{\num}[1]{\textbf{#1}}

% Numeric column that aligns on the decimal point.
\newcolumntype{R}{>{\raggedleft\arraybackslash}p{1.6cm}}

% Shared bar-chart styling. Use inside tikzpicture.
\pgfplotsset{
  barbase/.style={
    ybar,
    width=\linewidth,
    height=5.6cm,
    ymin=0,
    bar width=13pt,
    ymajorgrids=true,
    grid style={rule,dashed,very thin},
    axis line style={rule},
    tick align=outside,
    tick style={rule},
    legend style={draw=rule,font=\small,at={(0.5,1.04)},anchor=south,legend columns=-1},
    label style={font=\small},
    tick label style={font=\small},
    nodes near coords,
    nodes near coords style={font=\scriptsize,color=black!70},
  },
  stackbase/.style={
    ybar stacked,
    width=\linewidth,
    height=5.6cm,
    ymin=0,
    bar width=22pt,
    ymajorgrids=true,
    grid style={rule,dashed,very thin},
    axis line style={rule},
    legend style={draw=rule,font=\small,at={(1.02,1)},anchor=north west},
    label style={font=\small},
    tick label style={font=\small},
  },
}
 immediately after \documentclass is NOT needed;
% this file declares the class itself. Start your document with:
%   % latex-report skill: shared preamble
% Usage: \input{preamble} immediately after \documentclass is NOT needed;
% this file declares the class itself. Start your document with:
%   \input{preamble}
%   \begin{document}

\documentclass[11pt,a4paper]{article}

\usepackage[T1]{fontenc}
\usepackage[utf8]{inputenc}
\usepackage{lmodern}
\usepackage[margin=2.4cm,top=2.2cm,bottom=2.4cm]{geometry}
\usepackage{booktabs}
\usepackage{array}
\usepackage{longtable}
\usepackage{xcolor}
\usepackage{tcolorbox}
\usepackage{pgfplots}
\usepackage{caption}
\usepackage{titlesec}
\usepackage{enumitem}
\usepackage{fancyhdr}
\usepackage{microtype}
\usepackage[hidelinks]{hyperref}

\pgfplotsset{compat=1.18}
\usetikzlibrary{patterns}
\tcbuselibrary{skins}

% ---------------------------------------------------------------- palette
% Chosen to stay separable in greyscale and under deuteranopia.
\definecolor{palA}{HTML}{1F4E79}   % deep blue
\definecolor{palB}{HTML}{C55A11}   % burnt orange
\definecolor{palC}{HTML}{548235}   % green
\definecolor{palD}{HTML}{7030A0}   % purple
\definecolor{palE}{HTML}{808080}   % grey
\definecolor{rule}{HTML}{B0B0B0}
\definecolor{panelbg}{HTML}{F2F6FA}
\definecolor{panelln}{HTML}{1F4E79}

% ---------------------------------------------------------------- headings
\titleformat{\section}{\Large\bfseries\color{palA}}{\thesection}{0.6em}{}
\titleformat{\subsection}{\large\bfseries}{\thesubsection}{0.6em}{}
\titlespacing*{\section}{0pt}{1.4em}{0.6em}
\titlespacing*{\subsection}{0pt}{1.1em}{0.4em}

\captionsetup{font=small,labelfont={bf,color=palA},skip=6pt}
\setlist{itemsep=2pt,topsep=4pt,parsep=0pt}

\pagestyle{fancy}
\fancyhf{}
\renewcommand{\headrulewidth}{0.4pt}
\fancyfoot[C]{\small\thepage}

% ---------------------------------------------------------------- macros
% Boxed lead panel for the TL;DR.
\newcommand{\tldr}[1]{%
  \begin{tcolorbox}[colback=panelbg,colframe=panelln,boxrule=1pt,
    left=10pt,right=10pt,top=8pt,bottom=8pt,arc=2pt]
  #1
  \end{tcolorbox}\vspace{4pt}}

% Callout for a caveat or correction.
\newcommand{\caveat}[1]{%
  \begin{tcolorbox}[colback=white,colframe=palB,boxrule=0.8pt,
    left=9pt,right=9pt,top=6pt,bottom=6pt,arc=1pt]
  \small #1
  \end{tcolorbox}\vspace{3pt}}

\newcommand{\code}[1]{\texttt{\small #1}}
\newcommand{\num}[1]{\textbf{#1}}

% Numeric column that aligns on the decimal point.
\newcolumntype{R}{>{\raggedleft\arraybackslash}p{1.6cm}}

% Shared bar-chart styling. Use inside tikzpicture.
\pgfplotsset{
  barbase/.style={
    ybar,
    width=\linewidth,
    height=5.6cm,
    ymin=0,
    bar width=13pt,
    ymajorgrids=true,
    grid style={rule,dashed,very thin},
    axis line style={rule},
    tick align=outside,
    tick style={rule},
    legend style={draw=rule,font=\small,at={(0.5,1.04)},anchor=south,legend columns=-1},
    label style={font=\small},
    tick label style={font=\small},
    nodes near coords,
    nodes near coords style={font=\scriptsize,color=black!70},
  },
  stackbase/.style={
    ybar stacked,
    width=\linewidth,
    height=5.6cm,
    ymin=0,
    bar width=22pt,
    ymajorgrids=true,
    grid style={rule,dashed,very thin},
    axis line style={rule},
    legend style={draw=rule,font=\small,at={(1.02,1)},anchor=north west},
    label style={font=\small},
    tick label style={font=\small},
  },
}

%   \begin{document}

\documentclass[11pt,a4paper]{article}

\usepackage[T1]{fontenc}
\usepackage[utf8]{inputenc}
\usepackage{lmodern}
\usepackage[margin=2.4cm,top=2.2cm,bottom=2.4cm]{geometry}
\usepackage{booktabs}
\usepackage{array}
\usepackage{longtable}
\usepackage{xcolor}
\usepackage{tcolorbox}
\usepackage{pgfplots}
\usepackage{caption}
\usepackage{titlesec}
\usepackage{enumitem}
\usepackage{fancyhdr}
\usepackage{microtype}
\usepackage[hidelinks]{hyperref}

\pgfplotsset{compat=1.18}
\usetikzlibrary{patterns}
\tcbuselibrary{skins}

% ---------------------------------------------------------------- palette
% Chosen to stay separable in greyscale and under deuteranopia.
\definecolor{palA}{HTML}{1F4E79}   % deep blue
\definecolor{palB}{HTML}{C55A11}   % burnt orange
\definecolor{palC}{HTML}{548235}   % green
\definecolor{palD}{HTML}{7030A0}   % purple
\definecolor{palE}{HTML}{808080}   % grey
\definecolor{rule}{HTML}{B0B0B0}
\definecolor{panelbg}{HTML}{F2F6FA}
\definecolor{panelln}{HTML}{1F4E79}

% ---------------------------------------------------------------- headings
\titleformat{\section}{\Large\bfseries\color{palA}}{\thesection}{0.6em}{}
\titleformat{\subsection}{\large\bfseries}{\thesubsection}{0.6em}{}
\titlespacing*{\section}{0pt}{1.4em}{0.6em}
\titlespacing*{\subsection}{0pt}{1.1em}{0.4em}

\captionsetup{font=small,labelfont={bf,color=palA},skip=6pt}
\setlist{itemsep=2pt,topsep=4pt,parsep=0pt}

\pagestyle{fancy}
\fancyhf{}
\renewcommand{\headrulewidth}{0.4pt}
\fancyfoot[C]{\small\thepage}

% ---------------------------------------------------------------- macros
% Boxed lead panel for the TL;DR.
\newcommand{\tldr}[1]{%
  \begin{tcolorbox}[colback=panelbg,colframe=panelln,boxrule=1pt,
    left=10pt,right=10pt,top=8pt,bottom=8pt,arc=2pt]
  #1
  \end{tcolorbox}\vspace{4pt}}

% Callout for a caveat or correction.
\newcommand{\caveat}[1]{%
  \begin{tcolorbox}[colback=white,colframe=palB,boxrule=0.8pt,
    left=9pt,right=9pt,top=6pt,bottom=6pt,arc=1pt]
  \small #1
  \end{tcolorbox}\vspace{3pt}}

\newcommand{\code}[1]{\texttt{\small #1}}
\newcommand{\num}[1]{\textbf{#1}}

% Numeric column that aligns on the decimal point.
\newcolumntype{R}{>{\raggedleft\arraybackslash}p{1.6cm}}

% Shared bar-chart styling. Use inside tikzpicture.
\pgfplotsset{
  barbase/.style={
    ybar,
    width=\linewidth,
    height=5.6cm,
    ymin=0,
    bar width=13pt,
    ymajorgrids=true,
    grid style={rule,dashed,very thin},
    axis line style={rule},
    tick align=outside,
    tick style={rule},
    legend style={draw=rule,font=\small,at={(0.5,1.04)},anchor=south,legend columns=-1},
    label style={font=\small},
    tick label style={font=\small},
    nodes near coords,
    nodes near coords style={font=\scriptsize,color=black!70},
  },
  stackbase/.style={
    ybar stacked,
    width=\linewidth,
    height=5.6cm,
    ymin=0,
    bar width=22pt,
    ymajorgrids=true,
    grid style={rule,dashed,very thin},
    axis line style={rule},
    legend style={draw=rule,font=\small,at={(1.02,1)},anchor=north west},
    label style={font=\small},
    tick label style={font=\small},
  },
}

%   \begin{document}

\documentclass[11pt,a4paper]{article}

\usepackage[T1]{fontenc}
\usepackage[utf8]{inputenc}
\usepackage{lmodern}
\usepackage[margin=2.4cm,top=2.2cm,bottom=2.4cm]{geometry}
\usepackage{booktabs}
\usepackage{array}
\usepackage{longtable}
\usepackage{xcolor}
\usepackage{tcolorbox}
\usepackage{pgfplots}
\usepackage{caption}
\usepackage{titlesec}
\usepackage{enumitem}
\usepackage{fancyhdr}
\usepackage{microtype}
\usepackage[hidelinks]{hyperref}

\pgfplotsset{compat=1.18}
\usetikzlibrary{patterns}
\tcbuselibrary{skins}

% ---------------------------------------------------------------- palette
% Chosen to stay separable in greyscale and under deuteranopia.
\definecolor{palA}{HTML}{1F4E79}   % deep blue
\definecolor{palB}{HTML}{C55A11}   % burnt orange
\definecolor{palC}{HTML}{548235}   % green
\definecolor{palD}{HTML}{7030A0}   % purple
\definecolor{palE}{HTML}{808080}   % grey
\definecolor{rule}{HTML}{B0B0B0}
\definecolor{panelbg}{HTML}{F2F6FA}
\definecolor{panelln}{HTML}{1F4E79}

% ---------------------------------------------------------------- headings
\titleformat{\section}{\Large\bfseries\color{palA}}{\thesection}{0.6em}{}
\titleformat{\subsection}{\large\bfseries}{\thesubsection}{0.6em}{}
\titlespacing*{\section}{0pt}{1.4em}{0.6em}
\titlespacing*{\subsection}{0pt}{1.1em}{0.4em}

\captionsetup{font=small,labelfont={bf,color=palA},skip=6pt}
\setlist{itemsep=2pt,topsep=4pt,parsep=0pt}

\pagestyle{fancy}
\fancyhf{}
\renewcommand{\headrulewidth}{0.4pt}
\fancyfoot[C]{\small\thepage}

% ---------------------------------------------------------------- macros
% Boxed lead panel for the TL;DR.
\newcommand{\tldr}[1]{%
  \begin{tcolorbox}[colback=panelbg,colframe=panelln,boxrule=1pt,
    left=10pt,right=10pt,top=8pt,bottom=8pt,arc=2pt]
  #1
  \end{tcolorbox}\vspace{4pt}}

% Callout for a caveat or correction.
\newcommand{\caveat}[1]{%
  \begin{tcolorbox}[colback=white,colframe=palB,boxrule=0.8pt,
    left=9pt,right=9pt,top=6pt,bottom=6pt,arc=1pt]
  \small #1
  \end{tcolorbox}\vspace{3pt}}

\newcommand{\code}[1]{\texttt{\small #1}}
\newcommand{\num}[1]{\textbf{#1}}

% Numeric column that aligns on the decimal point.
\newcolumntype{R}{>{\raggedleft\arraybackslash}p{1.6cm}}

% Shared bar-chart styling. Use inside tikzpicture.
\pgfplotsset{
  barbase/.style={
    ybar,
    width=\linewidth,
    height=5.6cm,
    ymin=0,
    bar width=13pt,
    ymajorgrids=true,
    grid style={rule,dashed,very thin},
    axis line style={rule},
    tick align=outside,
    tick style={rule},
    legend style={draw=rule,font=\small,at={(0.5,1.04)},anchor=south,legend columns=-1},
    label style={font=\small},
    tick label style={font=\small},
    nodes near coords,
    nodes near coords style={font=\scriptsize,color=black!70},
  },
  stackbase/.style={
    ybar stacked,
    width=\linewidth,
    height=5.6cm,
    ymin=0,
    bar width=22pt,
    ymajorgrids=true,
    grid style={rule,dashed,very thin},
    axis line style={rule},
    legend style={draw=rule,font=\small,at={(1.02,1)},anchor=north west},
    label style={font=\small},
    tick label style={font=\small},
  },
}

\usepackage{amsmath,amssymb}
\usepackage{xurl}
\usepackage{multicol}
\usepackage{needspace}
\begin{document}
\emergencystretch=3em
\Urlmuskip=0mu plus 2mu

\newcommand{\unverified}{\textbf{UNVERIFIED}}
\newcommand{\opinion}{\textbf{Opinion.}}

\title{From language models to learned abstractions\\\large A research brief on pre-transformer lineages and the alternatives since 2017}
\author{Prepared for Andrew Oliver}
\date{11 October 2026}
\maketitle

\tldr{\textbf{Verdict.} Nothing published before or after 2017 has matched a transformer on general English while training on much less text. Each pre-2017 lineage solved one part of language and stalled on acquisition, coverage, search or memory cost. The post-2017 alternatives that reached production, state-space and linear-attention hybrids, cut inference cost on long inputs, still train on trillions of tokens, and all keep some attention for exact recall. The data-efficient results (MLC, ARC program search, AXIOM, QDisCoCirc, BabyLM) hold where structure was built in or generated. LeCun's AMI Labs has published vision and planning work and has not released a language model. \textbf{Opinion:} the most testable opening is a small learner meta-trained on generated compositional languages, with exact recall kept, compared against a transformer at matched compute and exposure past 100 million words.}

\section{How to read this brief}
This brief draws only on two research notes written in October 2026: one on lineages before the transformer and one on alternatives since 2017. Every bracketed number points to an entry in those notes, and the list at the end gives one link per work, numbered in order of first citation. Three markings carry over from the notes. \unverified{} marks a figure or attribution that the notes could not confirm against a primary source. \opinion{} marks judgment, either the notes' own or this brief's, and opinion sentences carry citations only for the evidence they lean on. Explanations of why a lineage lost adoption are synthesis by the research notes, which state that such an explanation is synthesis unless a source makes the comparison directly; those sentences carry no citation.

\section{The problem in numbers}
Children's spoken-word exposure by early childhood is in the tens of millions, with large differences between families and between counting methods \cite{pre47}. Hart and Risley extrapolated sampled hours to roughly 13, 26 and 45 million words by age four across three groups \cite{pre45}. Gilkerson and colleagues measured a mean of 12,709 adult words per recorded day over 3,213 all-day recordings, and the pre-2017 notes multiply that to about 18.6 million words over four years, a calculation that assumes a constant daily rate \cite{pre46}. Sperry, Sperry and Miller's study of 42 children in five communities did not support a uniform 30-million-word socioeconomic gap and argued that leaving out other caregivers and bystander speech undercounts some children's input \cite{pre47}. On the model side, GPT-3 trained on 300 billion tokens \cite{pre55}, Chinchilla on 1.4 trillion \cite{pre54}, Llama~3 405B on 15.6 trillion \cite{post84}, and Qwen3 on about 36 trillion, including trillions of synthetic tokens \cite{post85}.

\begin{table}[ht]
\centering\small
\begin{tabular}{p{3.0cm}p{3.4cm}p{7.6cm}}
\toprule
Evidence & Reported scale & What the number measures \\
\midrule
Hart and Risley & 13, 26 and 45 million words by age four & Extrapolation from about an hour a month of observation in 42 families \cite{pre45}. \\
Gilkerson et al. & 12,709 words per day; about 18.6 million by age four & Mean of 3,213 twelve-hour recordings; the four-year total is a constant-rate calculation in the notes \cite{pre46}. \\
GPT-3 & 300 billion tokens & Training tokens, with unequal repeated sampling of its datasets \cite{pre55}. \\
Chinchilla & 70B parameters, 1.4 trillion tokens & Beat 280B Gopher at the same training compute; about 20 tokens per parameter on its fitted frontier \cite{pre54}. \\
Llama 3, Qwen3 & 15.6 and about 36 trillion tokens & Frontier practice trains smaller models longer to cut inference cost \cite{post84,post85}. \\
BabyLM budgets & 10 and 100 million words & Challenge tracks set at or below a child's estimated input \cite{post81}. \\
\bottomrule
\end{tabular}
\caption{Exposure counts for children and training counts for models. Children's figures are spoken words and model figures are tokens.}
\label{tab:scale}
\end{table}

Table~\ref{tab:scale} sets the target. Dividing the model counts by Hart and Risley's range gives about 6,700 to 23,000 for GPT-3 and 31,000 to 108,000 for Chinchilla \cite{pre45,pre55,pre54}. Those ratios compare tokens with words, count repeated training passes again, and stop at age four while children keep learning. They measure the gap in raw exposure. The BabyLM organizers put the gap at three or four orders of magnitude \cite{post81}. Two further results frame it. Eight-month-old infants segmented an artificial speech stream from syllable statistics after two minutes of exposure, which isolates one small learning problem \cite{pre57}. At the other end, Muennighoff and colleagues found that up to four epochs of repeated data are nearly as good as unique data, after which ``the value of adding compute eventually decays to zero'' \cite{post86}. The scaling-law fits behind current practice describe transformer families in measured regimes \cite{pre53,pre54}. Whether a different learner needs trillions of tokens is an open question.

\section{What the transformer got right, and what it costs}
Vaswani and colleagues used attention for sequence transduction with no recurrent or convolutional sequence layers, computing $\mathrm{Attention}(Q,K,V)=\mathrm{softmax}(QK^\top/\sqrt{d_k})V$ with learned projections, multiple heads, positional representations, residual connections and feed-forward layers \cite{pre52}. The paper reports 28.4 BLEU on WMT14 English-to-German and 41.8 on English-to-French in its abstract and table, with 41.0 appearing in the text of the same version; the notes record that discrepancy \cite{pre52}. Its parts already existed. Bahdanau attention let a decoder consult every source state \cite{pre22}, convolutional sequence-to-sequence learning trained in parallel a month earlier \cite{pre51}, and backpropagation and recurrent networks had shown learned features and emergent lexical categories decades before \cite{pre58,pre59}. The combination gave position-parallel training under a causal mask and a one-step path between any two positions \cite{pre52}.

Attention also turned out to contain a learner. Induction heads copy what followed an earlier occurrence of the current token, and they appear at the same point in training as a jump in in-context learning \cite{post90}. A linear self-attention layer can implement one step of gradient descent on least squares over the prompt's examples, and self-attention models trained on regression tasks found weights matching that construction \cite{post91}; on LLaMA-7B with natural data, in-context learning and gradient descent respond differently to demonstration order, so the equivalence ``remains an open hypothesis'' \cite{post92}. The same softmax lookup is one update of a modern Hopfield network with capacity exponential in the dimension \cite{post21}. The post-2017 notes read the gated delta rule \cite{post8} as the same key-to-value fit made in a fixed-size state, one step of online regression per token, and Titans makes that step with momentum and weight decay \cite{post24}. Reasoning models then added variable inference computation: adaptive search or revision was more than four times as efficient as best-of-$N$, and in a FLOPs-matched comparison a small model with test-time compute beat a 14-times larger model on problems it could already partly solve \cite{post93}.

The costs are equally specific. Dense attention evaluates $n^2$ pairs at $O(n^2 d)$ cost, projections and feed-forward layers add roughly $O(nd^2)$, and a KV cache stores $O(nd)$ per layer during generation, so producing an $n$-token sequence still costs quadratic total attention work \cite{pre52}. Left-to-right factorization causes the reversal curse, which LLaDA's bidirectional training removed \cite{post101}, and pretraining runs consume 15 to 36 trillion tokens \cite{post84,post85}. Table~\ref{tab:keep} lists what the post-2017 notes find every serious replacement has had to keep or add back.

\begin{table}[ht]
\centering\small
\begin{tabular}{p{3.6cm}p{7.4cm}p{2.6cm}}
\toprule
Property & Evidence & Sources \\
\midrule
Parallel training & Every successful efficient model has a parallel or chunked training form, usually a structured matrix product & \cite{post1,post3,post5,post7} \\
Exact content recall & 82\% of a convolution-attention gap came from associative-recall tokens; fixed state bounds copying; MiniMax returned to full attention & \cite{post15,post16,post14} \\
In-context learning & Induction heads; gradient descent on the prompt in regression settings & \cite{post90,post91} \\
A generic objective & Next-token loss improves with data and compute along fitted curves & \cite{pre53,pre54} \\
Hardware fit & Dense matrix products; memory-bound linear kernels and settling methods lose much of their advantage & \cite{post14,post43} \\
Variable inference work & Adaptive test-time compute four times as efficient as best-of-$N$; o3 at 172 times the compute scored 87.5\% on ARC-AGI-1 & \cite{post93,post52} \\
\bottomrule
\end{tabular}
\caption{Properties a transformer replacement has had to keep, as summarized in the post-2017 notes.}
\label{tab:keep}
\end{table}

\section{Before 2017: what each lineage solved}
\label{sec:pre}
The earlier work produced distinct tools, each covering a bounded part of language. Each subsection below covers the core idea, the measured achievements, the drawbacks, and why the line stalled and what survived. Table~\ref{tab:pre} at the end of the section summarizes them.

\subsection{Formal grammar and symbolic NLP}
Formal grammar treats a language as the output of a finite rule set that generates unboundedly many structured expressions, with meaning computed from the parts according to the syntax. Chomsky's \emph{Syntactic Structures} (1957) formalized phrase structure and transformations \cite{pre1}, and Montague's \emph{Universal Grammar} (1970) mapped grammatical combination onto truth-conditional interpretation with intensional logic \cite{pre2}. Pullum and Scholz (2002) separate Chomsky's later poverty-of-the-stimulus argument from the 1957 book and dispute claims that decisive evidence for it exists \cite{pre3}. Engineering systems put the program to work. Woods's augmented transition network parser (1970) attached conditions and structure-building actions to grammar arcs \cite{pre61}; SHRDLU answered questions, executed commands and resolved some discourse ambiguity in a simulated blocks world \cite{pre44}; and MYCIN's rule-based medical advice received about 75\% expert approval in early studies, with later blinded comparisons showing that human experts also disagreed \cite{pre60}. The foundational works made recursion, constituent structure, quantifier scope and relations between sentences into explicit objects of analysis, and they report no held-out benchmark scores \cite{pre1,pre2}. The drawbacks were ambiguity, which multiplies analyses, and the cost of lexical coverage, pragmatics, exceptions and hand-maintained knowledge, with parser errors flowing downstream through pipelines. Statistical disambiguation and corpus-learned lexicons took over the coverage work, and neural models later removed most feature engineering. Parse structures, semantic types, scope distinctions, executable logical forms and constrained decoding survive.

\subsection{Statistical models: n-grams, HMMs and translation}
Statistical NLP estimates the probability of language events from counts and resolves ambiguity by choosing the likeliest analysis. Shannon (1948) built successive statistical approximations to English \cite{pre4}. Hidden Markov models factor a sequence through latent states and emissions, consolidated in Rabiner's 1989 tutorial \cite{pre7}, and the IBM translation models (1993) learned word alignments from bilingual corpora by expectation maximization \cite{pre6}. Kneser-Ney smoothing predicts unseen contexts from how many distinct predecessors a word follows; Kneser and Ney reported up to about 10\% lower perplexity and 5\% lower word error rate in trigram experiments \cite{pre5}. Conditional random fields (2001) allowed overlapping input features and globally normalized label sequences, which addressed label bias \cite{pre8}. Jelinek's own 2004 slides reproduce the remark ``Whenever I fire a linguist our system performance improves'' and date it to a December 1988 talk; the exact original wording is \unverified, and the anecdote supports no controlled comparison \cite{pre9}. Longer n-grams fragment counts across sparse contexts and grow storage, HMM state assumptions limit dependencies, CRFs need feature design and labeled data, and IBM translation needed parallel text and approximate decoding. Neural embeddings shared information across similar contexts, and encoder-decoder networks then displaced phrase-based translation on major benchmarks. Likelihood training, cross-entropy and perplexity, beam search, latent alignment and dynamic programming survive in every modern system. Shannon's theory explicitly brackets semantic meaning \cite{pre4}.

\subsection{Distributional semantics}
Distributional semantics infers word meaning from patterns of use, following Harris (1954) on substitutability \cite{pre10} and Firth's (1957) account of collocation \cite{pre11}. Latent semantic analysis (1990) factorized a term-document matrix with SVD into roughly 50 to 100 dimensions and retrieved documents through latent associations even when the query terms were absent \cite{pre12}. Word2vec (2013) learned vectors from 1.6 billion words at much lower cost than earlier neural embedding methods \cite{pre13}. The linguistic-regularities paper reported almost 40\% accuracy on its syntactic analogies \cite{pre14}, and GloVe's 300-dimensional vectors trained on 42 billion tokens reached 75.0\% overall analogy accuracy on a different test setup \cite{pre15}. Linzen (2016) measured how much of that arithmetic is neighborhood structure: without excluding the input words, the nearest neighbor of $a^*-a+b$ was $b$ itself in 93\% of tested cases, and reversing the analogy direction lowered average accuracy by 0.11 \cite{pre16}. A static vector mixes a word's senses, co-occurrence places antonyms beside synonyms, and local objectives say little about negation, role binding or quantifiers. Contextual neural representations took over downstream tasks because the same word could receive a different vector in each sentence. Embedding layers, cosine retrieval, dense indexing and similarity geometry survive. Word2vec's billion-word scale gives little support to a child-scale acquisition claim \cite{pre13}.

\subsection{Neural language models, recurrence and the attention bridge}
Neural language models share learned continuous features across contexts and predict the next word from a state. Bengio and colleagues' neural probabilistic language model (2003) learned word features jointly with a finite-context predictor; on Brown, the best n-gram baseline had 24\% higher test perplexity than the selected neural model \cite{pre17}. Elman (1990) had shown recurrent networks forming lexical categories from context \cite{pre59}, and Mikolov's recurrent LM (2010) reported around 50\% lower perplexity than a backoff model and around 18\% relative word-error reduction on WSJ with equal training data \cite{pre18}. Hochreiter and Schmidhuber's LSTM (1997) controlled error flow through memory cells and bridged delays over 1,000 steps on artificial tasks \cite{pre19}, and Cho and colleagues (2014) introduced the gated unit now called GRU \cite{pre20}. Sutskever, Vinyals and Le's sequence-to-sequence system scored 34.8 BLEU on WMT14 English-to-French against 33.3 for the cited phrase-based system, and 36.5 when reranking its hypotheses \cite{pre21}. Bahdanau, Cho and Bengio added learned soft alignment, so the decoder could consult different source states at each output step, and reached performance comparable to a phrase-based English-to-French system \cite{pre22}. Recurrence still trains one step at a time, a fixed-width state cannot hold arbitrary exact history in finite precision, and attending over every source position at each target step adds an $O(n_{\mathrm{src}}\,n_{\mathrm{tgt}})$ term. Transformers replaced recurrence with parallel attention and shorter paths between positions \cite{pre52}. Autoregressive factorization, embeddings, gates, encoder-decoder training and attention itself survive. Mikolov's results remain direct evidence that a learned representation can beat counting with less in-domain data, within a defined speech task \cite{pre18}.

\subsection{Recursive and tree-structured models}
Recursive networks compose a sentence along its parse tree, combining child vectors into a parent vector. Socher and colleagues' recursive tensor network and the Stanford Sentiment Treebank (2013) reached 85.4\% binary and 45.7\% fine-grained sentence sentiment accuracy \cite{pre23}. Tai, Socher and Manning's Tree-LSTM (2015) gated the memory flowing from each branch and reached 88.0\% binary and 51.0\% fine-grained accuracy with tuned GloVe vectors, plus a Pearson correlation of 0.8676 on semantic relatedness in its dependency form \cite{pre24}. Those results rest on extra information. An external parser supplied possibly wrong trees, phrase-level sentiment labels carry more supervision than sentence labels, and pretrained vectors bring outside text exposure. Tree-shaped computation also batches poorly because parents wait for their children. Sequence models simplified the pipeline, and attention learned cross-token interactions without an external tree; the notes treat this as an adoption explanation that leaves the value of syntax open. Recursive composition, graph message passing and phrase-level evaluation survive. The open experiment is whether learned structure lowers sample complexity when supervision and pretraining exposure are held equal.

\subsection{Compositional distributional semantics: DisCoCat}
DisCoCat lets grammar decide how word vectors combine. Coecke, Sadrzadeh and Clark (2010) joined Lambek pregroup grammar to vector spaces through shared category-theoretic structure. A transitive verb is a tensor in $N\otimes S\otimes N$, and the grammatical reduction $n\,(n^r s\, n^l)\,n \to s$ contracts subject and object against it, giving a sentence vector in one shared space $S$ whatever the sentence's grammatical shape; with Boolean scalars the model connects back to Montague semantics \cite{pre25}. Grefenstette and Sadrzadeh's 2011 implementation learned relational matrices from a corpus and reached Spearman correlations with human judgments of 0.17 on an intransitive test, matching a multiplicative baseline, and 0.21 on a transitive test \cite{pre26}. Those correlations leave most of the variation unexplained, and the work includes no general language-model comparison. A dense verb tensor has $d_N^2 d_S$ entries, which is $d^3$ when the widths match, and higher-order types grow further. Choosing the grammar, sparse words, negation and scope remain learning problems, and low-rank factorization saves storage by restricting the available operators. Grammar-led models needed more setup than broadly trained sequence networks, and the program continued as a research line and never became the dominant production architecture. Typed composition, tensor contractions and diagrammatic reasoning survive as the most direct mathematical statement of meaning as operators applied to arguments.

\subsection{Vector-symbolic architectures and hyperdimensional computing}
Vector-symbolic architectures bind a role to a filler and superpose the results, so structure can be stored in a distributed vector and approximately unpacked. Smolensky's tensor-product representation (1990) binds with outer products, $T=\sum_i f_i\otimes r_i$ \cite{pre27}. Plate's holographic reduced representations (1995) replaced the growing product with circular convolution, keeping a fixed width while encoding bindings, sequences and recursive frames, and demonstrated approximate reconstruction with a separate cleanup memory \cite{pre28}. Kanerva's sparse distributed memory (1988) and his 2009 tutorial, which illustrates the geometry with 10,000-bit vectors, gathered superposition, binding and permutation into one algebra \cite{pre30,pre29}. The sources contribute an algebra and a capacity analysis, with demonstrations of recall and structure and no broad language benchmark. Exact tensor products grow with structure depth, fixed-width versions accept interference and noisy reconstruction, cleanup needs a stored vocabulary searched at $O(Md)$ cost, and similarity between random codes carries no learned meaning. Elementwise binding costs $O(d)$ and FFT convolution $O(d\log d)$, so a fixed number of operations per word stays linear in sequence length before cleanup and capacity are counted. Gradient-trained representations gained broader task performance and software support, and these architectures stayed a smaller program. Role-filler separation, associative memory and approximate unbinding survive; the open question is how a system learns the codes and operators.

\subsection{Memory-augmented networks}
Memory-augmented networks separate a learned controller from addressable storage, so stored data can be reused without encoding everything in weights. The Neural Turing Machine (2014), trained on copy sequences of length 1 to 20, copied lengths 30 and 50 with few errors; its length-120 example had local errors and a duplicated item after the 128-location memory overwrote earlier writes. On associative recall it approached zero loss after about 30,000 episodes, while the comparison LSTM had not reached zero after a million \cite{pre31}. Memory Networks considered up to 14 million stored sentences and reported an eightfold speedup from clustered hashing at similar question-answering performance \cite{pre32}. End-to-end Memory Networks removed supporting-fact supervision and matched or slightly exceeded similarly sized RNN and LSTM language models on Penn Treebank and Text8, gaining from multiple hops \cite{pre33}. The Differentiable Neural Computer (2016) added allocation and temporal links and solved graph shortest paths and relational questions \cite{pre34}. Controllers must learn what to write, where and when, diffuse reads mix unrelated values, and a wrong write corrupts later reasoning. A soft scan costs $O(Md)$ per head per step and exact temporal linkage stores an $M\times M$ matrix, although the DNC manuscript also reports a sparse approximation with $O(M\log M)$ computation, $O(M)$ memory and no discernible loss on its tasks \cite{pre34}. The demonstrations stayed on selected tasks, and attention encoders and retrieval systems delivered the deployable parts. Content-based retrieval, multi-hop reads and controller-memory separation survive in retrieval-augmented systems.

\subsection{Bayesian program induction and the language of thought}
Bayesian program induction explains observations with a compact generative procedure and reuses its parts, scoring $P(\text{program}\mid D)\propto P(D\mid\text{program})\,P(\text{program})$. Fodor (1975) proposed a combinatorial language of thought \cite{pre35}. Tenenbaum (1998) explained generalization from a few positive examples with a size principle, $P(D\mid h)=|h|^{-m}$ for $m$ examples inside hypothesis $h$ \cite{pre36}, and Goodman, Tenenbaum and Gerstenberg (2014) treated concepts as stochastic functions \cite{pre38}. Lake, Salakhutdinov and Tenenbaum's Bayesian Program Learning (2015) represented handwritten characters as drawing programs and reached 3.3\% error on 20-way one-shot classification, against 4.5\% for people, 13.5\% for a deep convolutional network and 8.0\% for a task-optimized Siamese network. Removing compositionality raised its error to 14.0\%, and with only five background alphabets it scored 4.3\% and 4.0\% against 24.0\% and 22.3\% for the network \cite{pre37}. Piantadosi, Tenenbaum and Goodman modeled children's transition to recursive counting (2012) and used learning curves to compare candidate primitive sets (2016), favoring a richer set with first-order quantification \cite{pre39,pre40}. Program search is combinatorial, approximate posterior search can miss the right abstraction, and the hypothesis language decides which concepts are easy; BPL's stroke primitives encode knowledge about drawing. Dense networks offered batchable training with less domain setup, and the notes find no evidence that transformers disproved the approach. Hierarchical priors, program synthesis and reusable libraries survive, and BPL remains the clearest pre-2017 evidence that reusable structure buys few-shot learning, in a visual domain with supplied priors.

\subsection{The systematicity debate}
Fodor and Pylyshyn (1988) argued that a mind able to represent \emph{loves(John, Mary)} can also represent \emph{loves(Mary, John)}, and that this systematicity requires combinatorial syntax and semantics at the cognitive level even if neural networks implement it \cite{pre41}. It was an architectural argument without benchmark measurements. Its lasting effect was to make variable binding, productivity, role reversal and generalization to unseen combinations into explicit tests, and tensor-product binding and neural external memory were among the responses \cite{pre27,pre31}. The argument leaves open how primitives arise, how context changes meaning and how noise is handled, and a symbolic system can represent unlimited combinations while its search or acquisition stays impractical. Connectionist results showing learned distributed structure weakened the use of an unstructured associative network as a stand-in for all neural computation. The dispute became an empirical question about which architectures generalize under which distributions, which held-out accuracy alone never settled. The tests survive and return in SCAN, COGS and MLC (Section~\ref{sec:comp}).

\subsection{Smaller lines}
Four smaller lines complete the pre-2017 picture. Cognitive architectures coordinated rules, memory and goals: Soar (1987) cached subgoal results as productions through chunking \cite{pre42}, and the 2004 ACT-R synthesis related modules and buffers to cognitive and neural data \cite{pre43}, both with supplied task models. Grammar induction searched for the most compact grammar. Gold (1967) showed that the full class of regular languages cannot be identified from every arbitrary positive presentation under his criterion \cite{pre48}, Rissanen (1978) formalized minimum description length \cite{pre49}, and Stolcke and Omohundro (1994) induced finite-state and recursive grammars by Bayesian merging, with mixed results on natural corpora \cite{pre50}. Convolutional sequence-to-sequence learning, posted in May 2017, beat a deep LSTM on WMT14 translation with an order-of-magnitude speedup, so parallel training predates attention-only models \cite{pre51}. Latent Dirichlet allocation (2003) explained documents as mixtures of shared topics, a compact latent level above words that discards word order and role binding \cite{pre56}.

\begin{table}[ht]
\centering\small
\begin{tabular}{p{3.0cm}p{4.3cm}p{3.6cm}p{2.4cm}}
\toprule
Lineage & Best demonstrated contribution & Main drawback & Fate \\
\midrule
Formal grammar & Explicit recursion, scope and typed meaning & Coverage, ambiguity, hand-built knowledge & Components \\
Statistical NLP & Kneser-Ney: about 10\% lower perplexity & Sparse fixed context, feature work & Components \\
Distributional & GloVe: 75.0\% analogy accuracy & Similarity lacks composition & Embedding layers \\
Neural LM, RNN, LSTM & RNN mixture: about 50\% lower perplexity & Sequential training, finite state & Hybrids, attention \\
Tree models & Tree-LSTM: 88.0\% binary sentiment & Parser dependence, batching & Research line \\
DisCoCat & Typed composition; 0.21 correlation & Tensor cost, learned types & Research line \\
VSA and HDC & Binding algebra with fixed width & Interference, hand-made codes & Niche \\
Memory networks & NTM copy beyond trained lengths & Write control, access cost & Modern retrieval \\
Program induction & BPL: 3.3\% one-shot error & Search, domain priors & Formal tasks \\
Systematicity & Explicit tests of binding and role swap & No acquisition account & Benchmarks \\
\bottomrule
\end{tabular}
\caption{Pre-transformer lineages solved different subproblems. Figures come from Section~\ref{sec:pre}, where each carries its citation.}
\label{tab:pre}
\end{table}

\section{After 2017: alternatives and continuations}
\label{sec:post}

\subsection{Efficient sequence models and hybrids}
State-space models, linear attention and delta-rule memories replace the all-pairs comparison with a fixed-size state updated once per token. S4 (2021) averaged 80.48\% on Long Range Arena, where every earlier baseline was under 60\%, and solved Path-X at length 16,384 with 88\% \cite{post1}. Linear attention turned a causal transformer layer into an RNN with constant memory per step and ran up to 4,000 times faster on very long autoregressive sequences at similar quality \cite{post4}. RetNet reported 8.4 times faster decoding and 70\% less memory than a transformer with a KV cache, for a 7B model at 8k context \cite{post5}, and RWKV scaled a dense RNN to 14B parameters at reported parity \cite{post6}. Mamba made the state parameters depend on the input; Mamba-3B matched transformers twice its size with five times the generation throughput \cite{post2}. Dao and Gu showed that these families are views of one object, multiplication by a semiseparable matrix with an attention-like dual form, and a 2.7B Mamba-2 trained on 300B Pile tokens beat Pythia-6.9B trained on the same data \cite{post3}. The gated delta rule overwrites the value stored under the current key, the Widrow-Hoff rule applied as fast weights \cite{post8}; DeltaNet at 3B parameters and 1T tokens slightly underperformed a matched transformer \cite{post7}. NVIDIA's controlled study trained 8B models on up to 3.5T tokens. Pure SSMs lagged on five-shot MMLU, phonebook lookup and long-context reasoning, while a hybrid of 43\% Mamba-2, 7\% attention and 50\% MLP layers beat the transformer on all 12 standard tasks by 2.65 points on average \cite{post17}.

Two papers explain the lag. Zoology trained 17 models from 70M to 1.4B parameters on the Pile and traced 82\% of the perplexity gap between gated convolutions and attention to associative-recall tokens, which make up 6.4\% of Pile tokens. A 70M attention model beat a 1.4B Hyena model on those tokens (2.41 against 3.43 perplexity), and on multi-query associative recall the convolution models needed width growing linearly with sequence length while attention used a constant 64 \cite{post15}. Repeat After Me proved that a two-layer transformer can copy strings exponentially long in its size while any model with a fixed-size latent state is bounded by that state, and found pretrained transformers ahead of pretrained SSMs at copying and retrieval from context \cite{post16}. Production followed the same line. MiniMax-01, with 456B total and 45.9B active parameters, found that pure lightning attention failed needle-in-a-haystack retrieval and placed one softmax layer after every seven linear ones; its fitted scaling laws stayed in the same total-data regime \cite{post13}. MiniMax-M2 (October 2025) returned to full attention. Its pretraining lead wrote that the hybrid ``had clear deficits in complex, multi-hop reasoning tasks'' at scale, that linear attention is more sensitive to numerical precision, and that ``the race for efficient attention is a race to save compute'' \cite{post14}. AI21's Jamba interleaves Mamba and attention with 12B active of 52B parameters and an eight times smaller KV cache \cite{post9}, and Qwen3-Next repeats blocks of three Gated DeltaNet layers and one gated attention layer \cite{post10}. Alibaba reports that Qwen3.8-Flash-Next keeps that 3:1 ratio, costs about one ninth as much to train as Qwen3.7-Plus, and reaches a 7.6 times prefill kernel speedup at 1M tokens; these are vendor figures without independent replication, \unverified{} in the notes \cite{post11,post12}. Every production adopter kept some attention, and every model above trained on hundreds of billions to trillions of tokens, with linear-attention scaling fits in the same regime as softmax \cite{post13}. This line attacks compute. Its structured-matrix mathematics is reusable, and its recall results set a requirement: a replacement must still do exact content-addressed lookup over the context.

\subsection{Retrieval and memory}
Retrieval separates what a model knows from how it computes: facts sit in an external store that the network queries by similarity. Ramsauer and colleagues showed that softmax lookup is one update of a continuous modern Hopfield network that can store a number of patterns exponential in the dimension, so attention is already an associative memory over its context \cite{post21}. kNN-LM interpolated a trained model with nearest-neighbor search over (context, next token) pairs and cut Wikitext-103 perplexity to 15.79, 2.86 points better than the base model, with no extra training. A model trained on 100M tokens plus search over 3B tokens beat the same model trained on all 3B, and the authors concluded that ``learning similarity between sequences of text is easier than predicting the next word'' \cite{post18}. RETRO retrieved chunks from a 2-trillion-token database and matched GPT-3 and Jurassic-1 on the Pile with 25 times fewer parameters, with part of its perplexity gain coming from overlap between evaluation text and the database, which the paper reports by overlap level \cite{post19}. Memorizing Transformers improved steadily as memory grew to 262K tokens \cite{post20}. Meta's memory layers, scaled to 128B memory parameters and 1T tokens, beat dense models with more than twice the compute \cite{post23}, and Titans updates a neural memory by gradient descent at test time and scaled past 2M tokens of context \cite{post24}. The gains concentrate on facts. NVIDIA's reproduction of RETRO up to 9.5B parameters found it ``on par with GPT on other tasks'' \cite{post22}, and Meta's largest gains were factual \cite{post23}. Datastores are large, with one vector per training token for kNN-LM, the retriever's similarity is fixed by a separate encoder, and the base model still needs large-scale training to learn language. Retrieval-augmented generation is standard at inference, RETRO-style pretraining is niche, and learned memory as parameters is entering production as 51B parameters of n-gram embeddings in Qwen3.8-Flash-Next \cite{post12}.

\subsection{LeCun, JEPA, AMI Labs and the gap for language}
LeCun's 2022 position paper proposes an agent built around a world model trained by self-supervised prediction, with the Joint-Embedding Predictive Architecture (JEPA) at its center. Encoders map an observed part $x$ and a target $y$ to representations, and a predictor minimizes $D(s_y,\mathrm{Pred}(s_x,z))$, which reads as an energy that is low for compatible pairs \cite{post25}. Predicting in representation space lets the model drop detail it cannot predict, such as exact pixels or exact wording. The hazard is collapse to a constant encoder, which the 2022 paper handles with regularizers on latent information \cite{post25} and LeJEPA (2025) handles with a single regularizer, SIGReg, after proving that an isotropic Gaussian embedding minimizes downstream prediction risk \cite{post30}. I-JEPA trained a ViT-H/14 on ImageNet in under 1,200 GPU hours, more than ten times more efficient than pixel reconstruction with MAE \cite{post26}. V-JEPA reached 81.9\% on Kinetics-400 with a frozen backbone after training on 2M videos and lost less accuracy than reconstruction models when labels were scarce \cite{post27}. V-JEPA 2 pretrained on over 1M hours of video; its 300M-parameter action-conditioned predictor, post-trained on under 62 hours of robot video, planned pick-and-place zero-shot in two new labs at 16 seconds per action, against 4 minutes for the Cosmos world model \cite{post28}. LeWorldModel (March 2026) trained a 15M-parameter JEPA end to end from pixels on one GPU \cite{post31}. H-JEPA raised Visual AntMaze planning success from 18\% to 73\% with a three-level hierarchy \cite{post32}, and Meta FAIR's RoboJEPA, an 8B-parameter action-conditioned predictor trained on 23 robot datasets, has prediction error that follows a power law in compute \cite{post33}.

Language remains the gap in this program. LLM-JEPA (ICLR 2026) adds a JEPA term to next-token training and reports gains on NL-RX, GSM8K, Spider and RottenTomatoes across Llama~3, OpenELM, Gemma~2 and OLMo models. It needs two views of the same content, such as a description and its regex or SQL, and the authors call the absence of JEPA-style LLMs ``a testimony of the challenge in designing such objectives for language'' \cite{post29}. LeCun's slides argue that if each generated token leaves the set of correct answers with probability $e$, an $n$-token answer is correct with probability $(1-e)^n$ \cite{post34,post35}. \opinion{} The post-2017 notes judge that this assumes independent, unrecoverable errors, and that reasoning models which revise long chains of thought weaken both assumptions while leaving his point about learning from text alone standing. LeCun left Meta in November 2025 \cite{post36}. AMI Labs, based in Paris, announced a \$1.03B seed round at a \$3.5B valuation in March 2026, with LeCun as executive chairman and Alexandre LeBrun as CEO \cite{post36,post38}, and a stated aim of systems that understand the world, have persistent memory, and can reason and plan \cite{post37}. In October 2026 LeCun said AMI has 50 to 60 staff, is in a research phase and has no short-term product plan \cite{post39}. Its published work, LeWorldModel and H-JEPA, is robotics and planning \cite{post31,post32}, and AMI has not released a language model. The open question for language is what the two views of a sentence should be (paraphrase, translation, text with code, or text with perception) so that a JEPA can predict one from the other without reconstructing words. V-JEPA 2's million hours of video also show that the label efficiency appears downstream of large-scale pretraining \cite{post28}.

\subsection{Predictive coding, forward-forward and equilibrium propagation}
Local learning rules compute weight updates from quantities available at each connection, which could let analog or neuromorphic hardware learn without moving activations through memory. Predictive coding minimizes layerwise prediction errors, and Whittington and Bogacz (2017) showed that its weight change converges to the backpropagation gradient for certain parameters \cite{post40}. Equilibrium propagation lets an energy-based network settle, nudges it toward the target, and updates from the difference between the two equilibria, which equals the gradient as the nudge goes to zero \cite{post41}. Forward-forward trains each layer to give high ``goodness'' on real data and low goodness on negative data across two forward passes \cite{post42}. The published results are small-scale. Forward-forward reached 1.36\% test error on permutation-invariant MNIST against about 1.4\% for backpropagation, and 41\% to 46\% error on CIFAR-10 against 37\% to 39\% \cite{post42}. A symmetric three-phase estimate cut equilibrium propagation's CIFAR-10 error from 86.6\% to 11.7\%, within 0.6 points of backpropagation through time on the same network \cite{post43}. Predictive coding matched backpropagation up to about seven layers and degraded with depth, which the PCX benchmark traces to energy imbalances that shrink gradients exponentially \cite{post44}. Settling costs many iterations per example, so these methods run slower than backpropagation on GPUs, and forward-forward needs a source of negative data. None has trained a model at the scale of even a small language model \cite{post43,post44}. They approximate the same gradient on the same data, so their case rests on hardware cost. The tie they draw between learning and energy minimization links them to Hopfield memory and JEPA.

\subsection{Active inference}
Friston's free-energy principle holds that a system which keeps itself distinct from its environment behaves as if it minimizes variational free energy, an upper bound on surprise. Active inference adds action and chooses policies by minimizing expected free energy, which mixes goal-seeking with information-seeking \cite{post45}. The quantitative support comes from a dish of cultured rat neurons: a free-energy model fitted on sessions 1 to 10 predicted synaptic changes with under 4\% error and over 80\% of neuronal responses at session 100 \cite{post46}. VERSES AI's AXIOM, an object-centric agent with expanding Bayesian mixture models and no gradient descent, learned simplified pixel games within 10,000 interaction steps, beat deep RL baselines in that regime on the authors' own benchmark, and used an order of magnitude fewer parameters. The games use single-color sprites, and the paper reports no run on standard complex control tasks \cite{post47}. Bruineberg and colleagues argue that the literature conflates Markov blankets as an inference tool with a metaphysical boundary between agent and world \cite{post48}. A widely repeated criticism that the principle is unfalsifiable has an \unverified{} primary attribution in the notes. AXIOM's efficiency comes from hand-chosen priors about objects and piecewise-linear motion, close to the expert-system pattern Andrew wants to avoid. The notes found no active-inference language model that competes with standard ones. The idea worth keeping is an objective that makes a learner seek the data that would most reduce its uncertainty, which has no demonstration for language yet.

\needspace{6\baselineskip}
\subsection{Neurosymbolic methods: ARC, DreamCoder, AlphaGeometry, AlphaProof and TRM}
Neurosymbolic methods represent a solution as a program and use a network to guide the search for it, and a program found from a few examples generalizes exactly where its logic applies. DreamCoder (2021) alternates a wake phase that searches for programs with a sleep phase that refactors solutions into a growing library of abstractions and trains a neural search policy; learning library and policy together solved more problems faster across eight domains \cite{post49}. Chollet's ARC (2019) measures skill-acquisition efficiency with grid puzzles built on near-innate priors, and argues that ``unlimited priors or unlimited training data allow experimenters to `buy' arbitrary levels of skills'' \cite{post50}. In ARC Prize 2024 the eligible winner scored 53.5\% with an 8B model, augmentation and test-time training, and Greenblatt reached 42\% by having GPT-4o write and debug thousands of Python programs per task \cite{post51}. OpenAI's o3 scored 75.7\% on the semi-private set at \$26 per task and 87.5\% at 172 times the compute \cite{post52}. On ARC-AGI-2 in 2025, NVARC won with 24.03\% at about \$0.20 per task using synthetic data, test-time training and Tiny Recursive Model components \cite{post53}. Chollet's report named per-task refinement loops the theme of the year and said frontier reasoning remains ``fundamentally constrained to knowledge coverage'' \cite{post54}. The Tiny Recursive Model applies one 2-layer, 7M-parameter network recursively, trained on about 1,000 examples, and reached 44.6\% on ARC-AGI-1 and 7.8\% on ARC-AGI-2, against 15.8\% and 1.3\% for DeepSeek R1 \cite{post55}. ARC-AGI-3, launched in March 2026, is interactive, with results due 4 December 2026 \cite{post56,post57}. Third-party trackers list a top Kaggle score of 62.96 and 62.7 for GPT-6 Astra; these figures are \unverified{} against arcprize.org \cite{post58,post59}.

Formal mathematics shows the same recipe in its strongest setting. AlphaGeometry's language model, trained from scratch on 100 million synthetic theorems and proofs, proposes auxiliary constructions for a symbolic deduction engine; it solved 25 of 30 olympiad geometry problems against 10 for the previous best method, with symbolic deduction alone solving 14 \cite{post60}. AlphaProof searched for Lean proofs with reinforcement learning and, together with AlphaGeometry~2, scored 28 of 42 at IMO 2024. The problems were formalized by hand and each proof took two to three days of compute \cite{post61}. At IMO 2025 a Gemini Deep Think model scored 35 of 42 in natural language within the 4.5-hour limit, with no formal system \cite{post62}. The working pattern is a learned proposer, an exact checker and iteration. Its limits are domain-specific languages that do not transfer, checkers that exist only for formal domains, and a lead over plain language models that lasted about a year. The post-2017 notes call DreamCoder's library learning the closest existing mathematics to discovering the basis that ideas are written in, and it has not been made to work on natural language \cite{post49}.

\subsection{Compositional generalization: SCAN, COGS and MLC}
\label{sec:comp}
Compositional generalization asks whether a learner that knows the parts and the rules for combining them handles new combinations without new data. Formally, meaning is a homomorphism from syntax to semantics, $[\![f(a,b)]\!]=F([\![a]\!],[\![b]\!])$. SCAN (2018) maps commands such as ``jump twice and walk left'' to action sequences. Sequence-to-sequence RNNs generalized when test commands resembled training and failed when a primitive seen only in isolation had to be composed with others; the authors suggested that weak systematicity may explain part of neural networks' appetite for training data \cite{post63}. COGS (2020) held out combinations in a fragment of English semantic parsing, and transformers and LSTMs scored 96\% to 99\% in distribution and 16\% to 35\% on the generalization set, with 6 to 8 points of variance across seeds \cite{post64}. Meta-Learning for Compositionality (MLC, 2023) trained a standard transformer on a stream of episodes, each with a newly generated grammar for a pseudolanguage, a few study examples and queries. Its best run was fully systematic when choosing its top response, its sampled outputs were systematic 82.4\% of the time against 80.7\% for people, and it reproduced human error types; it also improved SCAN and COGS few-shot splits \cite{post65}. MLC's meta-training distribution comes from the authors' grammars, and the paper models adult compositional skill and leaves its acquisition aside \cite{post65}. Large pretrained models also improved on these benchmarks, which makes compositional skill hard to separate from coverage of similar patterns. What survives is the move of placing the inductive bias in the training distribution, minimizing $\mathbb{E}_{G}\,\mathbb{E}_{\text{examples}\sim G}\,[-\log p_\theta(\text{output}\mid\text{query},\text{study examples})]$ over generated grammars $G$, so the network is rewarded for inferring rules from a few examples.

\subsection{Category theory and quantum NLP}
Category theory describes systems by how their parts compose. In DisCoCat, sentence meaning is a functor from a grammar category to vector spaces, so composition holds by construction \cite{pre25}. Quantum circuits are a restricted family of tensor networks, so the same diagrams can run on quantum hardware, and the lambeq library turns parsed sentences into tensor networks or parameterized circuits \cite{post66}. QNLP in Practice (2023) ran the first NLP experiments on noisy quantum hardware with datasets above 100 sentences \cite{post67}. QDisCoCirc (2024) trained components classically on small texts and composed them into larger texts evaluated on Quantinuum's H1-1. Accuracy held on texts larger than those seen in training, while transformer, LSTM and GPT-4 baselines were near chance on the larger stories; the dataset is a handcrafted toy task about actors moving in directions \cite{post68}, and a companion paper proves its question answering BQP-hard \cite{post71}. Categorical deep learning proposes monad algebras as one theory covering architecture constraints and implementations \cite{post69}. Symbolica raised \$31M citing that paper \cite{post72} and announced an 85.28\% ARC-AGI-2 public-evaluation score in February 2026, whose connection to categorical methods is \unverified{} \cite{post73}. Geometric deep learning, which derives layers from a domain's symmetry group, is the part with broad results, in vision, chemistry and physics \cite{post70}. DisCoCat and QNLP results stay on hundreds of sentences or handcrafted stories, the models need a parser, high-order word tensors are expensive, and the 2024 categorical paper recovers existing architectures without a reported language-benchmark win. Language also lacks a clean symmetry group comparable to rotations for molecules. The functorial composition law survives as something a learner could keep while learning its types from data.

\subsection{Hyperdimensional computing and VSA since 2017}
Post-2017 work consolidated the VSA algebra of bundling, binding and permutation over near-orthogonal random vectors \cite{post75}, and its strongest results came in hardware and hybrids. IBM Research Zurich and ETH ran a complete HDC system on 760,000 phase-change memory devices and matched software accuracy on language identification, news classification and gesture recognition, because the codes tolerate device noise \cite{post76}. A neuro-vector-symbolic architecture reached 87.7\% on RAVEN and 88.1\% on I-RAVEN, with reasoning two orders of magnitude faster than earlier neuro-symbolic search \cite{post77}. A benchmark over 121 UCI datasets found HDC generalizing from few training instances and holding accuracy until about 40\% of memory bits were flipped, and also found that its gains ``could often be surpassed by off-the-shelf methods'' \cite{post78}. Practitioners report ``no clear guidelines'' for VSAs inside trained networks \cite{post79}, and a 2025 submission used VSA encodings to probe concepts in LLM embeddings at 0.94 binary accuracy \cite{post80}. A bundle stays recoverable only while its item count stays well below the dimension, noise grows with nesting depth, and codes are designed by hand. Attention's key-value lookup and the delta-rule update $S \leftarrow S + v k^\top$ are close relatives of binding and unbinding, which suggests that VSA also describes what trained networks already do. The missing piece is a data-efficient way to learn the symbols, which the neuro-vector-symbolic work supplies for vision only \cite{post77}.

\subsection{Data-efficient language learning: BabyLM and TinyStories}
BabyLM tests training at the budget a child hears: 100 million words in the Strict track and 10 million in Strict-Small \cite{post81}. In 2023 ELC-BERT won both tracks with LTG-BERT refinements and a learned mix of earlier layers, reaching a Strict aggregate of 0.74 against 0.60 for the OPT-125M baseline. It trained for over 450 epochs on Strict and over 2,000 on Strict-Small, and curriculum learning, the most common approach, was ``largely unsuccessful'' \cite{post81}. The 2024 winner, GPT-BERT, trained one transformer with a hybrid masked and causal objective \cite{post87}; no multimodal submission beat the baselines, and the organizers found ``a strong relationship between training FLOPs and average performance'' \cite{post82}. The 2025 edition added a compute cap and an interaction track. The MoEP mixture-of-experts model led Strict-Small with a macro average of 42.3 against 37.4 for the best baseline, no Strict submission beat the GPT-BERT baselines on aggregate, some models reached BLiMP grammaticality scores comparable to Llama 70B, and all stayed below Llama 70B and humans on GLUE \cite{post88}. TinyStories used GPT-3.5 and GPT-4 to write short stories in the vocabulary of a 3- to 4-year-old. Models under 10M parameters trained on it wrote fluent, consistent multi-paragraph stories, while 125M-parameter general models such as GPT-2 small rarely produced coherent text beyond a few words \cite{post89}. Grammar is learnable from 10 to 100 million words, while knowledge and reasoning lag \cite{post88}. The winning ideas were known transformer techniques run for hundreds of epochs, so these models save data and spend compute \cite{post81}, and TinyStories takes its knowledge from the teacher that wrote the corpus \cite{post89}. The notes found no model trained on 100 million words that is useful as a general language model. Teacher-generated data survives in frontier pretraining, including the trillions of synthetic tokens in Qwen3 \cite{post85}.

\subsection{KAN, liquid, spiking, diffusion and Large Concept Models}
\textbf{Kolmogorov-Arnold Networks} follow the theorem that writes a continuous multivariate function as sums of univariate compositions: each edge carries a learnable spline, and the network has no linear weights. Smaller KANs matched or beat larger MLPs at formula fitting and PDE solving \cite{post94}. A parameter- and FLOP-matched comparison found KANs ahead only on symbolic formula fitting, with MLPs ahead on machine learning, vision, NLP and audio; giving MLPs B-spline activations erased the formula advantage, and KANs forgot more in continual learning \cite{post95}. No language model of note uses KAN layers.

\textbf{Liquid networks} are continuous-time RNNs whose time constant depends on the input \cite{post96}. Liquid AI's shipped LFM2 came from a hardware-in-the-loop search that chose gated short convolutions for 10 of 16 blocks and grouped-query attention for the other 6, and reports twice Qwen3's decode and prefill speed on CPU \cite{post97,post98}. Its technical report states that adding linear-attention or state-space operators ``does not improve aggregate quality'' under on-device budgets once a few attention layers handle long-range retrieval \cite{post98}.

\textbf{Spiking language models} spend energy only when a neuron fires. SpikeGPT trained 45M to 260M-parameter spiking models with about five times less energy on neuromorphic hardware \cite{post99}. SpikingBrain converted Qwen2.5-7B into linear and hybrid spiking models with about 150B tokens of continual pretraining, which its authors put at about 2\% of typical pretraining data, and reported over 100 times faster time-to-first-token at 4M tokens. The starting checkpoint had already been pretrained on trillions of tokens, so the 2\% measures conversion cost \cite{post100}.

\textbf{Diffusion language models} start from a fully masked sequence and unmask tokens in parallel. LLaDA trained an 8B model from scratch on 2.3T tokens, beat LLaMA2 7B on nearly all of 15 tasks, matched LLaMA3 8B (trained on 15T tokens) on in-context learning, and removed the reversal curse. It cannot use a KV cache, and the token comparison mixes data quality with other differences \cite{post101}. Mercury Coder reached 1,109 and 737 tokens per second on H100s, up to ten times faster than speed-optimized frontier models at comparable coding quality \cite{post102}, and Google released an experimental Gemini Diffusion demo in May 2025 \cite{post105}.

\textbf{Large Concept Models} predict the next sentence embedding in Meta's frozen SONAR space and decode it back to text. A 7B diffusion variant trained on 2.7T tokens (Meta's publication page says 7.7T) showed zero-shot generalization across many languages in summarization \cite{post103}. The regression variant predicted the average of plausible next sentences, a point that ``may not correspond to any relevant point in SONAR space'' \cite{post103}, and a third-party review notes that much of the cross-lingual transfer comes from SONAR's alignment \cite{post104}. The run used trillions of tokens and did not beat token models on general tasks. Its averaging failure shows that one point in a fixed embedding space represents an idea with many valid expressions poorly \cite{post103}.

\begin{table}[ht]
\centering\scriptsize
\begin{tabular}{p{2.7cm}p{3.6cm}p{2.6cm}p{2.3cm}p{2.1cm}}
\toprule
Direction & Evidence & Main limitation & Attacks & Evidence level \\
\midrule
SSM and linear hybrids & Linear-time inference; hybrid beat transformer on 12 tasks & Recall; data unchanged & Compute & Production \\
Retrieval and memory & 100M trained tokens plus 3B searched beat 3B trained & Storage, search, facts only & Data partly, compute partly & Language \\
JEPA & Vision, video and planning; ten times cheaper than MAE & No general LM & Compute; data downstream & Vision; early language \\
Local learning rules & CIFAR-scale parity at shallow depth & No LM at any scale & Hardware compute & Early research \\
Active inference & Games within 10,000 steps & Hand-chosen object priors & Both, in games & Limited \\
Program synthesis, ARC & TRM 7M parameters: 44.6\% ARC-AGI-1 & Needs a formal checker & Both, in formal domains & Benchmarks \\
MLC & 82.4\% systematic against 80.7\% for people & Generator transfer & Data, synthetic & Academic \\
Category theory, QNLP & Larger toy stories where GPT-4 was near chance & Parser, scale, hardware & Data, toy domain & Early \\
HDC and VSA & 760,000-device in-memory system & Capacity, hand-made codes & Compute & Niche \\
BabyLM, TinyStories & BLiMP near Llama 70B from 100M words & Knowledge; teacher dependence & Data, narrow & Challenge results \\
Diffusion, LCM & Parallel decoding; concept-space transfer & Exact tokens; data scale & Compute; data uncertain & Research and demos \\
\bottomrule
\end{tabular}
\caption{Post-2017 directions target different resources and have evidence of different scope. Figures come from Section~\ref{sec:post}, where each carries its citation.}
\label{tab:post}
\end{table}

Table~\ref{tab:post} separates measured language evidence from vision, toy tasks and hardware claims. The post-2017 notes summarize the split: the compute-only directions are mature and in production, and the data-efficiency directions work where the target has formal structure that can be built in or generated \cite{post65,post68,post47,post70}. Figure~\ref{fig:map} places the directions on both axes.

\section{A map of compute and data evidence}
\begin{figure}[ht]
\centering
\begin{tikzpicture}[x=1cm,y=0.82cm,font=\footnotesize]
\fill[panelbg] (0,0) rectangle (13,7.6);
\draw[rule,dashed] (6.5,0) to (6.5,7.6);
\draw[rule,dashed] (0,3.8) to (13,3.8);
\draw[->,thick,palA] (0,0) to (13.3,0);
\draw[->,thick,palA] (0,0) to (0,7.9);
\node[anchor=north] at (6.5,-0.15) {Evidence of lower compute (training or inference), weak to strong};
\node[anchor=south,rotate=90] at (-0.15,3.8) {Evidence of lower training data, weak to strong};
\tikzset{lang/.style={circle,fill=palA,inner sep=2.2pt},toy/.style={circle,draw=palB,very thick,fill=white,inner sep=2.0pt}}
% compute and data both weak
\node[lang] (fa) at (0.6,0.6) {}; \node[anchor=west] at (fa.east) {Full-attention transformer};
\node[toy] (lr) at (2.6,1.5) {}; \node[anchor=west] at (lr.east) {Local learning rules};
\node[lang] (lcm) at (4.2,2.6) {}; \node[anchor=west] at (lcm.east) {LCM};
% compute strong, data weak
\node[lang] (ssm) at (10.6,0.7) {}; \node[anchor=east] at (ssm.west) {SSM and linear hybrids};
\node[lang] (kls) at (8.2,1.6) {}; \node[anchor=west] at (kls.east) {KAN, liquid, spiking};
\node[lang] (dif) at (9.6,2.9) {}; \node[anchor=west] at (dif.east) {Diffusion LMs};
% data strong, compute weak
\node[lang] (bab) at (0.6,6.0) {}; \node[anchor=west] at (bab.east) {BabyLM (many epochs)};
\node[toy] (mlc) at (3.6,6.9) {}; \node[anchor=west] at (mlc.east) {MLC};
\node[toy] (cat) at (1.4,4.6) {}; \node[anchor=west] at (cat.east) {DisCoCat, QNLP};
% both
\node[lang] (ret) at (6.5,3.8) {}; \node[anchor=north west] at (ret.south east) {Retrieval, memory (partial on both)};
\node[toy] (jep) at (10.0,4.8) {}; \node[anchor=west] at (jep.east) {JEPA};
\node[toy] (hdc) at (8.6,5.6) {}; \node[anchor=west] at (hdc.east) {HDC, VSA};
\node[toy] (ax) at (7.4,6.6) {}; \node[anchor=west] at (ax.east) {AXIOM};
\node[toy] (arc) at (9.9,7.1) {}; \node[anchor=west] at (arc.east) {TRM, ARC search};
% legend
\node[lang] (l1) at (0.4,-1.25) {}; \node[anchor=west] at (l1.east) {evidence in language};
\node[toy] (l2) at (4.0,-1.25) {}; \node[anchor=west] at (l2.east) {evidence only in vision, games, formal tasks or toy languages};
\end{tikzpicture}
\caption{Qualitative placement of each post-2017 direction by the evidence that it lowers compute (horizontal) and training data (vertical), following the ``attacks compute'' and ``attacks data'' table of the post-2017 notes. Positions are ordinal judgments with no measured coordinates. Retrieval sits at the center; every marker in the upper-right quadrant is open.}
\label{fig:map}
\end{figure}

Figure~\ref{fig:map} is an evidence classification. Retrieval moves data into a store and keeps the store's cost \cite{post18,post19}. JEPA's label efficiency follows large-scale pretraining \cite{post28}. BabyLM saves words and spends epochs \cite{post81}. Every direction in the upper-right quadrant has its evidence in a constrained domain, and no general language learner sits there with a demonstrated advantage over a transformer on both resources.

\section{Lessons and openings for new mathematics}
Read together, the pre-2017 subsections show the earlier lines stalling at one of four boundaries: representing ambiguity, acquiring useful primitives, preserving arbitrary context, or paying for search and external memory. The pre-2017 notes name three separable targets for a successor. Learned compositional operators should keep roles and support new combinations, with DisCoCat, tensor-product binding and systematicity supplying formalisms and tests \cite{pre25,pre27,pre41}. Reusable causal or procedural parts should cut the examples needed for a new concept, which BPL measured in a limited visual domain \cite{pre37}. Information access should have a controlled cost, which recurrence and hyperdimensional operations offer and memory networks show can slip back to quadratic storage \cite{pre34}. The post-2017 notes add that every data-efficient success shares one ingredient, a structural prior the learner did not discover from text: symmetry in geometric deep learning, objects in AXIOM, a domain-specific language in program synthesis, grammar types in DisCoCat, or a distribution of generated tasks in MLC and NVARC \cite{post70,post47,post49,post68,post65,post53}. Sutton's bitter lesson holds that ``general methods that leverage computation are ultimately the most effective'' and names symmetries and objects among the things to leave out \cite{post106}. AlphaProof's formal pipeline, matched a year later by a plain language model \cite{post62}, and MiniMax's return to full attention \cite{post14} fit that pattern. The pre-2017 notes add that a comparison has to count all prior data and supplied structure (pretrained embeddings, background alphabets, parse labels, hand-written primitives, generator programs, retrieval corpora) and record training FLOPs, search expansions, inference state and their growth with context.

\caveat{\opinion{} The five openings below follow the opinion section of the post-2017 notes. Each states what it would join and how it could fail.

\textbf{1. Learn the algebra.} DisCoCat has a composition law of the right shape, with meaning as a functor from grammar to vector spaces, and takes its parts from a hand-built grammar \cite{pre25}; DreamCoder learns a library of abstractions by compressing solutions \cite{post49}; categorical deep learning supplies vocabulary and no learning algorithm \cite{post69}. A model that learns the types and composition operations by compression, with composition enforced as a law, joins these lines. It could fail if any family of categories rich enough for English is too large to search from little data.

\textbf{2. Make the inductive bias a distribution over languages.} MLC and the ARC winners show a standard network becoming a systematic few-shot learner when trained on many generated tasks \cite{post65,post53}. Real text would then only identify which language this is. It could fail if English is atypical of the generator's distribution, so that real text still does most of the teaching.

\textbf{3. Represent an idea as a region or a distribution.} LCM's averaging failure and JEPA's collapse problem are the same problem \cite{post103,post25}; LeJEPA proves which embedding distribution is best \cite{post30}, and Hopfield energies and VSA superposition hold sets of meanings in one object \cite{post21,post75}. It could fail through collapse, or through decoding from a region to one sentence that is too slow or too vague to use.

\textbf{4. Keep exact recall and compress the rest.} The efficient-model literature converged on a compressed state learned by online regression plus a small amount of exact lookup \cite{post8,post13,post24}. It could fail if the exact store grows into the dominant cost, as full attention did.

\textbf{5. Train on two views of the same meaning.} LLM-JEPA needs paired views \cite{post29}, and GPT-BERT's two objectives on the same 100 million words beat either alone \cite{post87}. Paraphrase, translation, and text paired with code or perception are candidate pairs. It could fail if the alignment imports its knowledge from a large pretrained model, which moves the text cost into that model.}

\opinion{} The post-2017 notes propose the test that would separate an advance from a well-chosen prior: a learning curve of loss and novel-combination transfer against training tokens, compared with a transformer at matched compute and exposure and extended past 100 million words. The test set should include role reversal, negation, quantification, paraphrase and new combinations, and the report should include failure cases along with aggregate scores.

\section{Questions for the project}
The research notes leave four decisions open, and each one changes which mathematical proposal can be tested. First, the input: text only, or also images, speech, actions and interaction with a teacher, as in BabyLM's interaction track \cite{post88}. Second, which structures may be built in and which must be learned, given that every data-efficient success so far supplied a prior. Third, what observable behavior defines an \emph{idea}: a stable equivalence class of expressions, an executable procedure, a causal model, or a region in a learned space. Fourth, what result would convince Andrew: a held-out compositional benchmark, a data-versus-compute curve, transfer from a small corpus, or a working system on open-ended tasks.

\section{Glossary}
\begin{multicols}{2}
\small
\begin{description}[leftmargin=1em,itemsep=1pt]
\item[Attention] Content-dependent weighting of stored representations, computed from queries, keys and values.
\item[KV cache] The stored keys and values of earlier tokens, which grows with context during generation.
\item[Content-addressable recall] Retrieving information by matching its content to a query.
\item[Associative recall] Predicting a token by looking up a pairing seen earlier in the context; Zoology's MQAR task tests many such lookups at once.
\item[Induction head] A pair of attention heads that copies the continuation of an earlier occurrence of the current token.
\item[In-context learning] Adapting behavior from examples in a prompt without updating model parameters.
\item[State-space model] A sequence model that updates a compact internal state as inputs arrive.
\item[Delta rule] An update that overwrites the value stored under a key, equivalent to one step of online regression.
\item[JEPA] Joint-Embedding Predictive Architecture, which predicts a target representation from a context representation.
\item[Compositional generalization] Applying learned parts and rules to combinations absent from training.
\item[DisCoCat] Categorical compositional distributional semantics, in which grammar types determine tensor contractions of word meanings.
\item[VSA and HDC] Vector-symbolic architectures and hyperdimensional computing, which bind, superpose and permute high-dimensional vectors.
\item[ARC] Abstraction and Reasoning Corpus, grid tasks designed to test few-example generalization.
\item[BLiMP] A benchmark of minimal sentence pairs that scores grammatical knowledge.
\item[Compute-optimal] A model and data allocation that minimizes loss for a fixed compute budget in a measured training family.
\item[Perplexity] The exponential of average per-token cross-entropy; lower means better prediction.
\end{description}
\end{multicols}

\section*{References}
Each entry in this list comes from the numbered lists in the two research notes and is generated by \texttt{make\_refs.py}, which merges works that both notes cite. Items the notes flag carry the same flags in the text: the ARC-AGI-3 leaderboard figures and the link between Symbolica's ARC-AGI-2 result and categorical methods are \unverified{}, as are the primary attribution of the unfalsifiability criticism of the free-energy principle, Qwen3.8-Flash-Next's vendor cost and speed figures, and the exact original wording of Jelinek's 1988 remark. Entries the notes label secondary, third-party or vendor claims are cited only for what those sources report.

{\scriptsize
\renewcommand{\section}[2]{}% thebibliography's own heading is replaced by the one above
\makeatletter\renewcommand{\@openbib@code}{\setlength{\itemsep}{0pt}\setlength{\parsep}{0pt}}\makeatother
\urlstyle{same}\renewcommand{\UrlFont}{\scriptsize}\raggedright
\setlength{\columnsep}{1.2em}
\begin{multicols}{2}
% Generated by make_refs.py. Do not edit by hand.
% Sources: research-pre2017.md (61 entries), research-post2017.md (106 entries).
% Aliases (same normalized URL or title; citations rewritten to the primary key):
%   post74 -> pre25
%   post83 -> pre54
\begin{thebibliography}{999}
\bibitem{pre47} Sperry, D. E.; Sperry, L. L.; Miller, P. J. \emph{Reexamining the Verbal Environments of Children From Different Socioeconomic Backgrounds}. Child Development, 2018 online; 2019 issue. \url{https://onlinelibrary.wiley.com/doi/abs/10.1111/cdev.13072}
\bibitem{pre45} Hart, B.; Risley, T. R. \emph{The Early Catastrophe: The 30 Million Word Gap by Age 3}. American Educator; author excerpt from their 1995 book, \emph{Meaningful Differences in the Everyday Experience of Young American Children}, 2003. \url{https://www.aft.org/ae/spring2003/hart_risley}
\bibitem{pre46} Gilkerson, J. et al. \emph{Mapping the Early Language Environment Using All-Day Recordings and Automated Analysis}. American Journal of Speech-Language Pathology, 2017. \url{https://doi.org/10.1044/2016_AJSLP-15-0169}
\bibitem{pre55} Brown, T. B. et al. \emph{Language Models are Few-Shot Learners}. 2020. \url{https://arxiv.org/abs/2005.14165v4}
\bibitem{pre54} Hoffmann, J. et al. \emph{Training Compute-Optimal Large Language Models}. 2022. \url{https://arxiv.org/abs/2203.15556}
\bibitem{post84} Llama Team (Meta). \emph{The Llama 3 Herd of Models}. 2024. \url{https://arxiv.org/abs/2407.21783}
\bibitem{post85} Qwen Team. \emph{Qwen3 Technical Report}. 2025. \url{https://arxiv.org/abs/2505.09388}
\bibitem{post81} Warstadt et al. \emph{Findings of the BabyLM Challenge: Sample-Efficient Pretraining on Developmentally Plausible Corpora}. CoNLL 2023. \url{https://aclanthology.org/2023.conll-babylm.1.pdf}
\bibitem{pre57} Saffran, J. R.; Aslin, R. N.; Newport, E. L. \emph{Statistical Learning by 8-Month-Old Infants}. Science, 1996. \url{https://www.science.org/doi/10.1126/science.274.5294.1926}
\bibitem{post86} Muennighoff et al. \emph{Scaling Data-Constrained Language Models}. JMLR 26, 2025. \url{https://www.jmlr.org/beta/papers/v26/24-1000.html}
\bibitem{pre53} Kaplan, J. et al. \emph{Scaling Laws for Neural Language Models}. 2020. \url{https://arxiv.org/abs/2001.08361}
\bibitem{pre52} Vaswani, A. et al. \emph{Attention Is All You Need}. 2017. \url{https://arxiv.org/abs/1706.03762}
\bibitem{pre22} Bahdanau, D.; Cho, K.; Bengio, Y. \emph{Neural Machine Translation by Jointly Learning to Align and Translate}. 2014 preprint; ICLR 2015. \url{https://arxiv.org/abs/1409.0473}
\bibitem{pre51} Gehring, J. et al. \emph{Convolutional Sequence to Sequence Learning}. ICML, 2017. \url{https://proceedings.mlr.press/v70/gehring17a.html}
\bibitem{pre58} Rumelhart, D. E.; Hinton, G. E.; Williams, R. J. \emph{Learning Representations by Back-Propagating Errors}. Nature, 1986. \url{https://www.nature.com/articles/323533a0}
\bibitem{pre59} Elman, J. L. \emph{Finding Structure in Time}. Cognitive Science, 1990. \url{https://onlinelibrary.wiley.com/doi/10.1207/s15516709cog1402_1}
\bibitem{post90} Olsson et al. (Anthropic). \emph{In-context Learning and Induction Heads}. 2022. \url{https://arxiv.org/abs/2209.11895}
\bibitem{post91} von Oswald et al. \emph{Transformers Learn In-Context by Gradient Descent}. ICML 2023. \url{https://proceedings.mlr.press/v202/von-oswald23a.html}
\bibitem{post92} Shen, Mishra, Khashabi. \emph{Position: Do pretrained Transformers Learn In-Context by Gradient Descent?} ICML 2024. \url{https://proceedings.mlr.press/v235/shen24d.html}
\bibitem{post21} Ramsauer et al. \emph{Hopfield Networks is All You Need}. 2020 (ICLR 2021). \url{https://arxiv.org/abs/2008.02217}
\bibitem{post8} Yang, Kautz, Hatamizadeh. \emph{Gated Delta Networks: Improving Mamba2 with Delta Rule}. 2024. \url{https://arxiv.org/abs/2412.06464}
\bibitem{post24} Behrouz, Zhong, Mirrokni (Google). \emph{Titans: Learning to Memorize at Test Time}. 2025. \url{https://arxiv.org/abs/2501.00663}
\bibitem{post93} Snell et al. \emph{Scaling LLM Test-Time Compute Optimally can be More Effective than Scaling Model Parameters}. 2024. \url{https://arxiv.org/abs/2408.03314}
\bibitem{post101} Nie et al. \emph{Large Language Diffusion Models (LLaDA)}. NeurIPS 2025. \url{https://proceedings.neurips.cc/paper_files/paper/2025/file/48b383b24230e0e6e649d9c98dae4d8c-Paper-Conference.pdf}
\bibitem{post1} Gu, Goel, R\'{e}. \emph{Efficiently Modeling Long Sequences with Structured State Spaces (S4)}. 2021. \url{https://arxiv.org/abs/2111.00396}
\bibitem{post3} Dao, Gu. \emph{Transformers are SSMs: Generalized Models and Efficient Algorithms Through Structured State Space Duality (Mamba-2)}. ICML 2024. \url{https://arxiv.org/abs/2405.21060}
\bibitem{post5} Sun et al. \emph{Retentive Network: A Successor to Transformer for Large Language Models}. 2023. \url{https://arxiv.org/abs/2307.08621}
\bibitem{post7} Yang et al. \emph{Parallelizing Linear Transformers with the Delta Rule over Sequence Length}. NeurIPS 2024. \url{https://proceedings.nips.cc/paper_files/paper/2024/file/d13a3eae72366e61dfdc7eea82eeb685-Paper-Conference.pdf}
\bibitem{post15} Arora et al. \emph{Zoology: Measuring and Improving Recall in Efficient Language Models}. ICLR 2024. \url{https://proceedings.iclr.cc/paper_files/paper/2024/file/448fc91f669c15d10364ee01d512cc10-Paper-Conference.pdf}
\bibitem{post16} Jelassi et al. \emph{Repeat After Me: Transformers are Better than State Space Models at Copying}. ICML 2024. \url{https://proceedings.mlr.press/v235/jelassi24a.html}
\bibitem{post14} MiniMax (M2 pretraining lead). \emph{Why Did M2 End Up as a Full Attention Model?} 29 October 2025. \url{https://www.minimax.io/news/why-did-m2-end-up-as-a-full-attention-model}
\bibitem{post43} Laborieux et al. \emph{Scaling Equilibrium Propagation to Deep ConvNets by Drastically Reducing Its Gradient Estimator Bias}. Frontiers in Neuroscience, 2021. \url{https://www.frontiersin.org/journals/neuroscience/articles/10.3389/fnins.2021.633674/full}
\bibitem{post52} Chollet. \emph{OpenAI o3 Breakthrough High Score on ARC-AGI-Pub}. ARC Prize blog, 20 December 2024. \url{https://arcprize.org/blog/oai-o3-pub-breakthrough}
\bibitem{pre1} Chomsky, N. \emph{Syntactic Structures}. Mouton, 1957. \url{https://babel.ucsc.edu/~hank/synstruct.pdf}
\bibitem{pre2} Montague, R. \emph{Universal Grammar}. Theoria, 1970. \url{https://onlinelibrary.wiley.com/doi/10.1111/j.1755-2567.1970.tb00434.x}
\bibitem{pre3} Pullum, G. K.; Scholz, B. C. \emph{Empirical Assessment of Stimulus Poverty}. The Linguistic Review, 2002. \url{https://sites.socsci.uci.edu/~lpearl/courses/readings/PullumScholz2002_PoSReassessment.pdf}
\bibitem{pre61} Woods, W. A. \emph{Transition Network Grammars for Natural Language Analysis}. Communications of the ACM, 1970. \url{https://dl.acm.org/doi/abs/10.1145/355598.362773}
\bibitem{pre44} Winograd, T. \emph{Procedures as a Representation for Data in a Computer Program for Understanding Natural Language}. MIT AITR-235, 1971 report; 1970 dissertation origin. \url{https://dspace.mit.edu/entities/publication/3406b826-5a41-4b72-9430-7876c0d5405e}
\bibitem{pre60} Buchanan, B. G.; Shortliffe, E. H. (eds.). \emph{Rule-Based Expert Systems: The MYCIN Experiments of the Stanford Heuristic Programming Project}. Addison-Wesley, 1984. \url{https://www.shortliffe.net/Buchanan-Shortliffe-1984/Chapter-01.pdf}
\bibitem{pre4} Shannon, C. E. \emph{A Mathematical Theory of Communication}. Bell System Technical Journal, 1948. \url{https://www.princeton.edu/~wbialek/rome/refs/shannon_48.pdf}
\bibitem{pre7} Rabiner, L. R. \emph{A Tutorial on Hidden Markov Models and Selected Applications in Speech Recognition}. Proceedings of the IEEE, 1989. \url{https://web.mit.edu/6.435/www/Rabiner89.pdf}
\bibitem{pre6} Brown, P. F. et al. \emph{The Mathematics of Statistical Machine Translation: Parameter Estimation}. Computational Linguistics, 1993. \url{https://aclanthology.org/www.mt-archive.info/CL-1993-Brown.pdf}
\bibitem{pre5} Kneser, R.; Ney, H. \emph{Improved Backing-off for M-gram Language Modeling}. ICASSP, 1995. \url{https://www-i6.informatik.rwth-aachen.de/publications/download/951/KneserReinhardNeyHermann--ImprovedBacking-offforM-gramLanguageModeling--1995.pdf}
\bibitem{pre8} Lafferty, J.; McCallum, A.; Pereira, F. \emph{Conditional Random Fields: Probabilistic Models for Segmenting and Labeling Sequence Data}. ICML, 2001. \url{https://courses.ischool.berkeley.edu/i290-dm/s11/SECURE/crf.pdf}
\bibitem{pre9} Jelinek, F. \emph{Some of my Best Friends are Linguists}. LREC slides, 2004. \url{http://www.lrec-conf.org/lrec2004/doc/jelinek.pdf}
\bibitem{pre10} Harris, Z. S. \emph{Distributional Structure}. WORD, 1954. \url{https://www.tandfonline.com/doi/pdf/10.1080/00437956.1954.11659520}
\bibitem{pre11} Firth, J. R. \emph{A Synopsis of Linguistic Theory, 1930-1955}. In \emph{Studies in Linguistic Analysis}, 1957. \url{https://languagelog.ldc.upenn.edu/myl/Firth1957.pdf}
\bibitem{pre12} Deerwester, S. et al. \emph{Indexing by Latent Semantic Analysis}. Journal of the American Society for Information Science, 1990. \url{https://web.stanford.edu/class/linguist289/lsi.pdf}
\bibitem{pre13} Mikolov, T. et al. \emph{Efficient Estimation of Word Representations in Vector Space}. 2013. \url{https://arxiv.org/abs/1301.3781}
\bibitem{pre14} Mikolov, T.; Yih, W.-t.; Zweig, G. \emph{Linguistic Regularities in Continuous Space Word Representations}. NAACL, 2013. \url{https://aclanthology.org/N13-1090.pdf}
\bibitem{pre15} Pennington, J.; Socher, R.; Manning, C. D. \emph{GloVe: Global Vectors for Word Representation}. EMNLP, 2014. \url{https://aclanthology.org/D14-1162/}
\bibitem{pre16} Linzen, T. \emph{Issues in Evaluating Semantic Spaces Using Word Analogies}. RepEval, 2016. \url{https://aclanthology.org/W16-2503/}
\bibitem{pre17} Bengio, Y. et al. \emph{A Neural Probabilistic Language Model}. JMLR, 2003. \url{https://jmlr.org/papers/v3/bengio03a.html}
\bibitem{pre18} Mikolov, T. et al. \emph{Recurrent Neural Network Based Language Model}. Interspeech, 2010. \url{https://www.fit.vut.cz/research/group/speech/public/publi/2010/mikolov_interspeech2010_IS100722.pdf}
\bibitem{pre19} Hochreiter, S.; Schmidhuber, J. \emph{Long Short-Term Memory}. Neural Computation, 1997. \url{https://dl.acm.org/doi/10.1162/neco.1997.9.8.1735}
\bibitem{pre20} Cho, K. et al. \emph{Learning Phrase Representations Using RNN Encoder-Decoder for Statistical Machine Translation}. 2014. \url{https://arxiv.org/abs/1406.1078}
\bibitem{pre21} Sutskever, I.; Vinyals, O.; Le, Q. V. \emph{Sequence to Sequence Learning with Neural Networks}. 2014. \url{https://arxiv.org/abs/1409.3215}
\bibitem{pre23} Socher, R. et al. \emph{Recursive Deep Models for Semantic Compositionality Over a Sentiment Treebank}. EMNLP, 2013. \url{https://aclanthology.org/D13-1170/}
\bibitem{pre24} Tai, K. S.; Socher, R.; Manning, C. D. \emph{Improved Semantic Representations From Tree-Structured Long Short-Term Memory Networks}. ACL-IJCNLP, 2015. \url{https://aclanthology.org/P15-1150/}
\bibitem{pre25} Coecke, B.; Sadrzadeh, M.; Clark, S. \emph{Mathematical Foundations for a Compositional Distributional Model of Meaning}. 2010. \url{https://arxiv.org/abs/1003.4394}
\bibitem{pre26} Grefenstette, E.; Sadrzadeh, M. \emph{Experimental Support for a Categorical Compositional Distributional Model of Meaning}. EMNLP, 2011. \url{https://aclanthology.org/D11-1129/}
\bibitem{pre27} Smolensky, P. \emph{Tensor Product Variable Binding and the Representation of Symbolic Structures in Connectionist Systems}. Artificial Intelligence, 1990. \url{https://doi.org/10.1016/0004-3702(90)90007-M}
\bibitem{pre28} Plate, T. A. \emph{Holographic Reduced Representations}. IEEE Transactions on Neural Networks, 1995. \url{https://doi.org/10.1109/72.377968}
\bibitem{pre30} Kanerva, P. \emph{Sparse Distributed Memory}. MIT Press, 1988. \url{https://mitpress.mit.edu/9780262111324/sparse-distributed-memory/}
\bibitem{pre29} Kanerva, P. \emph{Hyperdimensional Computing: An Introduction to Computing in Distributed Representation with High-Dimensional Random Vectors}. Cognitive Computation, 2009. \url{https://doi.org/10.1007/s12559-009-9009-8}
\bibitem{pre31} Graves, A.; Wayne, G.; Danihelka, I. \emph{Neural Turing Machines}. 2014. \url{https://arxiv.org/html/1410.5401v2}
\bibitem{pre32} Weston, J.; Chopra, S.; Bordes, A. \emph{Memory Networks}. 2014 preprint; ICLR 2015. \url{https://ar5iv.labs.arxiv.org/html/1410.3916}
\bibitem{pre33} Sukhbaatar, S. et al. \emph{End-To-End Memory Networks}. NeurIPS, 2015. \url{https://proceedings.neurips.cc/paper_files/paper/2015/file/8fb21ee7a2207526da55a679f0332de2-Paper.pdf}
\bibitem{pre34} Graves, A. et al. \emph{Hybrid Computing Using a Neural Network with Dynamic External Memory}. Nature, 2016. \url{https://www.nature.com/articles/nature20101}
\bibitem{pre35} Fodor, J. A. \emph{The Language of Thought}. Thomas Y. Crowell; Harvard University Press edition, 1975. \url{https://danielwharris.com/teaching/38080/readings/FodorLOT.pdf}
\bibitem{pre36} Tenenbaum, J. B. \emph{Bayesian Modeling of Human Concept Learning}. NIPS 1998. \url{https://proceedings.neurips.cc/paper/1998/file/d010396ca8abf6ead8cacc2c2f2f26c7-Paper.pdf}
\bibitem{pre38} Goodman, N. D.; Tenenbaum, J. B.; Gerstenberg, T. \emph{Concepts in a Probabilistic Language of Thought}. CBMM Memo 010, 2014 memo. \url{https://cbmm.mit.edu/publications/concepts-probabilistic-language-thought}
\bibitem{pre37} Lake, B. M.; Salakhutdinov, R.; Tenenbaum, J. B. \emph{Human-level Concept Learning through Probabilistic Program Induction}. Science, 2015. \url{https://www.science.org/doi/10.1126/science.aab3050}
\bibitem{pre39} Piantadosi, S. T.; Tenenbaum, J. B.; Goodman, N. D. \emph{Bootstrapping in a Language of Thought: A Formal Model of Numerical Concept Learning}. Cognition, 2012. \url{https://doi.org/10.1016/j.cognition.2011.11.005}
\bibitem{pre40} Piantadosi, S. T.; Tenenbaum, J. B.; Goodman, N. D. \emph{The Logical Primitives of Thought: Empirical Foundations for Compositional Cognitive Models}. Psychological Review, 2016. \url{https://doi.org/10.1037/a0039980}
\bibitem{pre41} Fodor, J. A.; Pylyshyn, Z. W. \emph{Connectionism and Cognitive Architecture: A Critical Analysis}. Cognition, 1988. \url{https://www.sciencedirect.com/science/article/pii/0010027788900315}
\bibitem{pre42} Laird, J. E.; Newell, A.; Rosenbloom, P. S. \emph{Soar: An Architecture for General Intelligence}. Artificial Intelligence, 1987. \url{https://raw.githubusercontent.com/SoarGroup/website-downloads/main/pubs/the_soar_papers_vol_1/15.pdf}
\bibitem{pre43} Anderson, J. R. et al. \emph{An Integrated Theory of the Mind}. Psychological Review, 2004. \url{https://doi.org/10.1037/0033-295X.111.4.1036}
\bibitem{pre48} Gold, E. M. \emph{Language Identification in the Limit}. Information and Control, 1967. \url{https://doi.org/10.1016/S0019-9958(67)91165-5}
\bibitem{pre49} Rissanen, J. \emph{Modeling by Shortest Data Description}. Automatica, 1978. \url{https://www.sciencedirect.com/science/article/pii/0005109878900055}
\bibitem{pre50} Stolcke, A.; Omohundro, S. M. \emph{Inducing Probabilistic Grammars by Bayesian Model Merging}. 1994. \url{https://arxiv.org/abs/cmp-lg/9409010}
\bibitem{pre56} Blei, D. M.; Ng, A. Y.; Jordan, M. I. \emph{Latent Dirichlet Allocation}. JMLR, 2003. \url{https://www.jmlr.org/papers/v3/blei03a.html}
\bibitem{post4} Katharopoulos et al. \emph{Transformers are RNNs: Fast Autoregressive Transformers with Linear Attention}. ICML 2020. \url{https://proceedings.mlr.press/v119/katharopoulos20a.html}
\bibitem{post6} Peng et al. \emph{RWKV: Reinventing RNNs for the Transformer Era}. 2023. \url{https://arxiv.org/abs/2305.13048}
\bibitem{post2} Gu, Dao. \emph{Mamba: Linear-Time Sequence Modeling with Selective State Spaces}. 2023. \url{https://arxiv.org/abs/2312.00752}
\bibitem{post17} Waleffe et al. (NVIDIA). \emph{An Empirical Study of Mamba-based Language Models}. 2024. \url{https://research.nvidia.com/publication/2024-06_empirical-study-mamba-based-language-models}
\bibitem{post13} MiniMax. \emph{MiniMax-01: Scaling Foundation Models with Lightning Attention}. 2025. \url{https://arxiv.org/abs/2501.08313}
\bibitem{post9} Lieber, Lenz et al. (AI21). \emph{Jamba: A Hybrid Transformer-Mamba Language Model}. 2024. \url{https://arxiv.org/abs/2403.19887}
\bibitem{post10} Qwen team. \emph{Qwen3-Next-80B-A3B-Instruct model card}. 2025. \url{https://huggingface.co/Qwen/Qwen3-Next-80B-A3B-Instruct}
\bibitem{post11} Qwen team. \emph{Qwen3.8-Flash-Next model card and repository}. August 2026. \url{https://huggingface.co/Qwen/Qwen3.8-Flash-Next}
\bibitem{post12} Alibaba Cloud Community. \emph{Qwen3.8-Flash-Next: A New Architecture, Towards Ultimate Cost-Efficiency}. 27 August 2026. \url{https://www.alibabacloud.com/blog/qwen-3-8-flash-next-a-new-architecture-towards-ultimate-cost-efficiency_603501}
\bibitem{post18} Khandelwal et al. \emph{Generalization through Memorization: Nearest Neighbor Language Models}. ICLR 2020. \url{https://arxiv.org/abs/1911.00172}
\bibitem{post19} Borgeaud et al. (DeepMind). \emph{Improving Language Models by Retrieving from Trillions of Tokens (RETRO)}. ICML 2022. \url{https://proceedings.mlr.press/v162/borgeaud22a.html}
\bibitem{post20} Wu et al. \emph{Memorizing Transformers}. ICLR 2022. \url{https://research.google/pubs/memorizing-transformers/}
\bibitem{post23} Berges et al. (Meta). \emph{Memory Layers at Scale}. December 2024. \url{https://ai.meta.com/research/publications/memory-layers-at-scale/}
\bibitem{post22} Wang, Ping et al. (NVIDIA). \emph{Shall We Pretrain Autoregressive Language Models with Retrieval? A Comprehensive Study}. EMNLP 2023. \url{https://aclanthology.org/2023.emnlp-main.482/}
\bibitem{post25} LeCun. \emph{A Path Towards Autonomous Machine Intelligence}. 2022. \url{https://openreview.net/pdf?id=BZ5a1r-kVsf}
\bibitem{post30} Balestriero, LeCun. \emph{LeJEPA: Provable and Scalable Self-Supervised Learning Without the Heuristics}. November 2025. \url{https://arxiv.org/abs/2511.08544}
\bibitem{post26} Assran et al. \emph{Self-Supervised Learning from Images with a Joint-Embedding Predictive Architecture (I-JEPA)}. CVPR 2023. \url{https://arxiv.org/abs/2301.08243}
\bibitem{post27} Bardes et al. \emph{Revisiting Feature Prediction for Learning Visual Representations from Video (V-JEPA)}. 2024. \url{https://arxiv.org/abs/2404.08471}
\bibitem{post28} Assran et al. (Meta). \emph{V-JEPA 2: Self-Supervised Video Models Enable Understanding, Prediction and Planning}. 2025. \url{https://arxiv.org/abs/2506.09985}
\bibitem{post31} Maes et al. \emph{LeWorldModel: Stable End-to-End Joint-Embedding Predictive Architecture from Pixels}. March 2026. \url{https://arxiv.org/abs/2603.19312}
\bibitem{post32} Zhang et al. \emph{H-JEPA: End-to-End Learning of Hierarchical World Models for Visual Planning}. arXiv 2610.06805, October 2026. \url{https://h-jepa.com/}
\bibitem{post33} Meta FAIR (correspondence Zholus, Bohg). \emph{RoboJEPA: Scaling Robotic Latent World Models}. October 2026. \url{https://arxiv.org/html/2610.10515}
\bibitem{post29} Huang, LeCun, Balestriero. \emph{LLM-JEPA: Large Language Models Meet Joint Embedding Predictive Architectures}. ICLR 2026. \url{https://proceedings.iclr.cc/paper_files/paper/2026/hash/ac818ec2da1c976b267de32f59709c69-Abstract-Conference.html}
\bibitem{post34} LeCun. \emph{Objective-Driven AI (slides, Harvard CMSA, 28 March 2024)}. 2024. \url{https://cmsa.fas.harvard.edu/media/lecun-20240328-harvard_reduced.pdf}
\bibitem{post35} LeCun. \emph{LinkedIn post on autoregressive LLMs as exponentially diverging processes, 27 March 2023}. 2023. \url{https://www.linkedin.com/posts/yann-lecun_i-have-claimed-that-auto-regressive-llms-activity-7045908925660950528-hJGk}
\bibitem{post36} Zeff. \emph{Yann LeCun Raises \$1 Billion to Build AI That Understands the Physical World}. WIRED, 10 March 2026 (secondary). \url{https://www.wired.com/story/yann-lecun-raises-dollar1-billion-to-build-ai-that-understands-the-physical-world/}
\bibitem{post38} Abdullahi. \emph{Yann LeCun, Meta's Former AI Chief, Launches \$1B Startup Focused on `World Models'}. eWeek, 10 March 2026 (secondary). \url{https://www.eweek.com/news/meta-ai-yann-lecun-ami-world-models-startup/}
\bibitem{post37} LeCun. \emph{Unveiling our new startup Advanced Machine Intelligence (AMI Labs)}. LinkedIn, 10 March 2026. \url{https://www.linkedin.com/posts/yann-lecun_unveiling-our-new-startup-advanced-machine-activity-7437016101080883200-0Mr9}
\bibitem{post39} LeCun. \emph{Interview at AIM 2026, transcript (secondary)}. October 2026. \url{https://ceointerviews.ai/interview/2961732/}
\bibitem{post40} Whittington, Bogacz. \emph{An Approximation of the Error Backpropagation Algorithm in a Predictive Coding Network with Local Hebbian Synaptic Plasticity}. Neural Computation, 2017. \url{https://exa.ai/library/publication/1q76nrfw1js}
\bibitem{post41} Scellier, Bengio. \emph{Equilibrium Propagation: Bridging the Gap between Energy-Based Models and Backpropagation}. Frontiers in Computational Neuroscience, 2017. \url{https://www.frontiersin.org/journals/computational-neuroscience/articles/10.3389/fncom.2017.00024/full}
\bibitem{post42} Hinton. \emph{The Forward-Forward Algorithm: Some Preliminary Investigations}. 2022. \url{https://arxiv.org/abs/2212.13345}
\bibitem{post44} Pinchetti et al. \emph{Benchmarking Predictive Coding Networks, Made Simple}. 2024. \url{https://arxiv.org/abs/2407.01163}
\bibitem{post45} Friston et al. \emph{The free energy principle made simpler but not too simple}. Physics Reports, 2023. \url{https://www.sciencedirect.com/science/article/pii/S037015732300203X}
\bibitem{post46} Isomura et al. \emph{Experimental validation of the free-energy principle with in vitro neural networks}. Nature Communications, 2023. \url{https://www.nature.com/articles/s41467-023-40141-z}
\bibitem{post47} Heins et al. (VERSES AI). \emph{AXIOM: Learning to Play Games in Minutes with Expanding Object-Centric Models}. 2025. \url{https://arxiv.org/abs/2505.24784}
\bibitem{post48} Bruineberg et al. \emph{The Emperor's New Markov Blankets}. Behavioral and Brain Sciences 45, 2022. \url{https://www.cambridge.org/core/journals/behavioral-and-brain-sciences/article/abs/emperors-new-markov-blankets/715C589A73DDF861DCF8997271DE0B8C}
\bibitem{post49} Ellis et al. \emph{DreamCoder: Bootstrapping Inductive Program Synthesis with Wake-Sleep Library Learning}. PLDI 2021. \url{https://dl.acm.org/doi/10.1145/3453483.3454080}
\bibitem{post50} Chollet. \emph{On the Measure of Intelligence}. 2019. \url{https://arxiv.org/abs/1911.01547}
\bibitem{post51} Chollet et al. \emph{ARC Prize 2024: Technical Report}. 2024. \url{https://arxiv.org/abs/2412.04604}
\bibitem{post53} Knoop. \emph{ARC Prize 2025 Results and Analysis}. ARC Prize blog, 5 December 2025. \url{https://arcprize.org/blog/arc-prize-2025-results-analysis}
\bibitem{post54} Chollet et al. \emph{ARC Prize 2025: Technical Report}. January 2026. \url{https://arxiv.org/abs/2601.10904}
\bibitem{post55} Jolicoeur-Martineau. \emph{Less is More: Recursive Reasoning with Tiny Networks}. October 2025. \url{https://arxiv.org/abs/2510.04871}
\bibitem{post56} ARC Prize. \emph{ARC Prize 2026 competition page}. 2026. \url{https://arcprize.org/competitions/2026}
\bibitem{post57} Kamradt. \emph{Announcing ARC-AGI-3}. ARC Prize blog, 25 March 2026. \url{https://arcprize.org/blog/arc-agi-3-launch}
\bibitem{post58} ARC Explainer (third-party tracker). \emph{ARC-AGI-3 leaderboard snapshot, 11 October 2026}. 2026. \url{https://arc.markbarney.net/arc-leaderboards}
\bibitem{post59} LLMLearner (third-party tracker). \emph{ARC-AGI-3 model rankings, October 2026}. 2026. \url{https://llmlearner.com/rankings/arc-agi-3}
\bibitem{post60} Trinh et al. \emph{Solving olympiad geometry without human demonstrations (AlphaGeometry)}. Nature, 2024. \url{https://www.nature.com/articles/s41586-023-06747-5}
\bibitem{post61} Hubert et al. (Google DeepMind). \emph{Olympiad-level formal mathematical reasoning with reinforcement learning (AlphaProof)}. Nature, November 2025. \url{https://www.nature.com/articles/s41586-025-09833-y}
\bibitem{post62} Luong, Lockhart. \emph{Advanced version of Gemini with Deep Think officially achieves gold-medal standard at the International Mathematical Olympiad}. Google DeepMind blog, 21 July 2025. \url{https://deepmind.google/blog/advanced-version-of-gemini-with-deep-think-officially-achieves-gold-medal-standard-at-the-international-mathematical-olympiad/}
\bibitem{post63} Lake, Baroni. \emph{Generalization without Systematicity: On the Compositional Skills of Sequence-to-Sequence Recurrent Networks (SCAN)}. ICML 2018. \url{https://proceedings.mlr.press/v80/lake18a.html}
\bibitem{post64} Kim, Linzen. \emph{COGS: A Compositional Generalization Challenge Based on Semantic Interpretation}. EMNLP 2020. \url{https://aclanthology.org/2020.emnlp-main.731/}
\bibitem{post65} Lake, Baroni. \emph{Human-like systematic generalization through a meta-learning neural network}. Nature, 2023. \url{https://www.nature.com/articles/s41586-023-06668-3}
\bibitem{post66} Kartsaklis et al. \emph{lambeq: An Efficient High-Level Python Library for Quantum NLP}. 2021. \url{https://arxiv.org/abs/2110.04236}
\bibitem{post67} Lorenz et al. \emph{QNLP in Practice: Running Compositional Models of Meaning on a Quantum Computer}. JAIR, 2023. \url{https://www.jair.org/index.php/jair/article/view/14329}
\bibitem{post68} Duneau et al. \emph{Scalable and interpretable quantum natural language processing: an implementation on trapped ions}. 2024. \url{https://arxiv.org/abs/2409.08777}
\bibitem{post71} Laakkonen, Meichanetzidis, Coecke. \emph{Quantum Algorithms for Compositional Text Processing}. 2024. \url{https://arxiv.org/abs/2408.06061}
\bibitem{post69} Gavranovi\'{c} et al. \emph{Position: Categorical Deep Learning is an Algebraic Theory of All Architectures}. ICML 2024. \url{https://proceedings.mlr.press/v235/gavranovic24a.html}
\bibitem{post72} Tong. \emph{Symbolica raises \$31 mln to develop AI systems to compete with OpenAI}. Reuters, 9 April 2024 (secondary). \url{https://www.reuters.com/technology/symbolica-raises-31-mln-develop-ai-systems-compete-with-openai-2024-04-09/}
\bibitem{post73} Symbolica AI. \emph{``Today we set a new ARC-AGI-2 public eval SOTA: 85.28\%.''} LinkedIn, 12 February 2026 (vendor claim). \url{https://www.linkedin.com/posts/symbolica-ai_today-we-set-a-new-arc-agi-2-public-eval-activity-7427546074229137408-28WU}
\bibitem{post70} Bronstein et al. \emph{Geometric Deep Learning: Grids, Groups, Graphs, Geodesics, and Gauges}. 2021. \url{https://arxiv.org/abs/2104.13478}
\bibitem{post75} Kleyko et al. \emph{A Survey on Hyperdimensional Computing aka Vector Symbolic Architectures, Part I}. ACM Computing Surveys, 2022. \url{https://dl.acm.org/doi/10.1145/3538531}
\bibitem{post76} Karunaratne et al. \emph{In-memory hyperdimensional computing}. Nature Electronics, 2020. \url{https://www.nature.com/articles/s41928-020-0410-3}
\bibitem{post77} Hersche et al. \emph{A neuro-vector-symbolic architecture for solving Raven's progressive matrices}. Nature Machine Intelligence, 2023. \url{https://www.nature.com/articles/s42256-023-00630-8}
\bibitem{post78} Verg\'{e}s et al. \emph{Classification using hyperdimensional computing: a review with comparative analysis}. Artificial Intelligence Review, 2025. \url{https://link.springer.com/article/10.1007/s10462-025-11181-2}
\bibitem{post79} Carzaniga et al. \emph{Practical Lessons on Vector-Symbolic Architectures in Deep Learning-Inspired Environments}. NeSy 2025. \url{https://proceedings.mlr.press/v284/carzaniga25a.html}
\bibitem{post80} Authors not recorded in the retrieved copy. \emph{Hyperdimensional Probe: Decoding LLM Representations via Vector Symbolic Architectures}. OpenReview submission, 2025. \url{https://openreview.net/pdf?id=WxM7lIoGBb}
\bibitem{post87} Charpentier, Samuel. \emph{GPT or BERT: why not both?} BabyLM 2024. \url{https://aclanthology.org/2024.conll-babylm.24/}
\bibitem{post82} Hu et al. \emph{Findings of the Second BabyLM Challenge}. CoNLL 2024. \url{https://aclanthology.org/2024.conll-babylm.1/}
\bibitem{post88} Charpentier et al. \emph{Findings of the Third BabyLM Challenge: Accelerating Language Modeling Research with Cognitively Plausible Data}. 2025. \url{https://aclanthology.org/2025.babylm-main.28/}
\bibitem{post89} Eldan, Li. \emph{TinyStories: How Small Can Language Models Be and Still Speak Coherent English?} 2023. \url{https://arxiv.org/abs/2305.07759}
\bibitem{post94} Liu et al. \emph{KAN: Kolmogorov-Arnold Networks}. 2024. \url{https://arxiv.org/abs/2404.19756}
\bibitem{post95} Yu, Yu, Wang. \emph{KAN or MLP: A Fairer Comparison}. 2024. \url{https://arxiv.org/abs/2407.16674}
\bibitem{post96} Hasani et al. \emph{Liquid Time-constant Networks}. AAAI 2021. \url{https://ojs.aaai.org/index.php/AAAI/article/view/16936}
\bibitem{post97} Liquid AI. \emph{Introducing LFM2: The Fastest On-Device Foundation Models on the Market}. 10 July 2025. \url{https://www.liquid.ai/blog/liquid-foundation-models-v2-our-second-series-of-generative-ai-models}
\bibitem{post98} Liquid AI. \emph{LFM2 Technical Report}. November 2025. \url{https://arxiv.org/abs/2511.23404}
\bibitem{post99} Zhu, Zhao, Eshraghian. \emph{SpikeGPT: Generative Pre-trained Language Model with Spiking Neural Networks}. 2023. \url{https://arxiv.org/abs/2302.13939}
\bibitem{post100} BICLab. \emph{SpikingBrain Technical Report: Spiking Brain-inspired Large Models}. September 2025. \url{https://arxiv.org/abs/2509.05276}
\bibitem{post102} Inception Labs. \emph{Mercury: Ultra-Fast Language Models Based on Diffusion}. 2025. \url{https://arxiv.org/abs/2506.17298}
\bibitem{post105} Google. \emph{Gemini Diffusion: Google DeepMind's experimental research model}. 20 May 2025. \url{https://blog.google/innovation-and-ai/models-and-research/google-deepmind/gemini-diffusion/}
\bibitem{post103} LCM team (Meta). \emph{Large Concept Models: Language Modeling in a Sentence Representation Space}. December 2024. \url{https://arxiv.org/abs/2412.08821}
\bibitem{post104} Pith (third-party review). \emph{Review of Large Concept Models}. n.d. \url{https://pith.science/paper/2412.08821}
\bibitem{post106} Sutton. \emph{The Bitter Lesson}. 13 March 2019. \url{https://web.archive.org/web/20251006184925/http:/www.incompleteideas.net/IncIdeas/BitterLesson.html}
\end{thebibliography}

\end{multicols}
}

\end{document}
